% Options for packages loaded elsewhere
\PassOptionsToPackage{unicode}{hyperref}
\PassOptionsToPackage{hyphens}{url}
\documentclass[
]{article}
\usepackage{xcolor}
\usepackage{amsmath,amssymb}
\setcounter{secnumdepth}{-\maxdimen} % remove section numbering
\usepackage{iftex}
\ifPDFTeX
  \usepackage[T1]{fontenc}
  \usepackage[utf8]{inputenc}
  \usepackage{textcomp} % provide euro and other symbols
\else % if luatex or xetex
  \usepackage{unicode-math} % this also loads fontspec
  \defaultfontfeatures{Scale=MatchLowercase}
  \defaultfontfeatures[\rmfamily]{Ligatures=TeX,Scale=1}
\fi
\usepackage{lmodern}
\ifPDFTeX\else
  % xetex/luatex font selection
\fi
% Use upquote if available, for straight quotes in verbatim environments
\IfFileExists{upquote.sty}{\usepackage{upquote}}{}
\IfFileExists{microtype.sty}{% use microtype if available
  \usepackage[]{microtype}
  \UseMicrotypeSet[protrusion]{basicmath} % disable protrusion for tt fonts
}{}
\makeatletter
\@ifundefined{KOMAClassName}{% if non-KOMA class
  \IfFileExists{parskip.sty}{%
    \usepackage{parskip}
  }{% else
    \setlength{\parindent}{0pt}
    \setlength{\parskip}{6pt plus 2pt minus 1pt}}
}{% if KOMA class
  \KOMAoptions{parskip=half}}
\makeatother
\usepackage{color}
\usepackage{fancyvrb}
\newcommand{\VerbBar}{|}
\newcommand{\VERB}{\Verb[commandchars=\\\{\}]}
\DefineVerbatimEnvironment{Highlighting}{Verbatim}{commandchars=\\\{\}}
% Add ',fontsize=\small' for more characters per line
\newenvironment{Shaded}{}{}
\newcommand{\AlertTok}[1]{\textcolor[rgb]{1.00,0.00,0.00}{\textbf{#1}}}
\newcommand{\AnnotationTok}[1]{\textcolor[rgb]{0.38,0.63,0.69}{\textbf{\textit{#1}}}}
\newcommand{\AttributeTok}[1]{\textcolor[rgb]{0.49,0.56,0.16}{#1}}
\newcommand{\BaseNTok}[1]{\textcolor[rgb]{0.25,0.63,0.44}{#1}}
\newcommand{\BuiltInTok}[1]{\textcolor[rgb]{0.00,0.50,0.00}{#1}}
\newcommand{\CharTok}[1]{\textcolor[rgb]{0.25,0.44,0.63}{#1}}
\newcommand{\CommentTok}[1]{\textcolor[rgb]{0.38,0.63,0.69}{\textit{#1}}}
\newcommand{\CommentVarTok}[1]{\textcolor[rgb]{0.38,0.63,0.69}{\textbf{\textit{#1}}}}
\newcommand{\ConstantTok}[1]{\textcolor[rgb]{0.53,0.00,0.00}{#1}}
\newcommand{\ControlFlowTok}[1]{\textcolor[rgb]{0.00,0.44,0.13}{\textbf{#1}}}
\newcommand{\DataTypeTok}[1]{\textcolor[rgb]{0.56,0.13,0.00}{#1}}
\newcommand{\DecValTok}[1]{\textcolor[rgb]{0.25,0.63,0.44}{#1}}
\newcommand{\DocumentationTok}[1]{\textcolor[rgb]{0.73,0.13,0.13}{\textit{#1}}}
\newcommand{\ErrorTok}[1]{\textcolor[rgb]{1.00,0.00,0.00}{\textbf{#1}}}
\newcommand{\ExtensionTok}[1]{#1}
\newcommand{\FloatTok}[1]{\textcolor[rgb]{0.25,0.63,0.44}{#1}}
\newcommand{\FunctionTok}[1]{\textcolor[rgb]{0.02,0.16,0.49}{#1}}
\newcommand{\ImportTok}[1]{\textcolor[rgb]{0.00,0.50,0.00}{\textbf{#1}}}
\newcommand{\InformationTok}[1]{\textcolor[rgb]{0.38,0.63,0.69}{\textbf{\textit{#1}}}}
\newcommand{\KeywordTok}[1]{\textcolor[rgb]{0.00,0.44,0.13}{\textbf{#1}}}
\newcommand{\NormalTok}[1]{#1}
\newcommand{\OperatorTok}[1]{\textcolor[rgb]{0.40,0.40,0.40}{#1}}
\newcommand{\OtherTok}[1]{\textcolor[rgb]{0.00,0.44,0.13}{#1}}
\newcommand{\PreprocessorTok}[1]{\textcolor[rgb]{0.74,0.48,0.00}{#1}}
\newcommand{\RegionMarkerTok}[1]{#1}
\newcommand{\SpecialCharTok}[1]{\textcolor[rgb]{0.25,0.44,0.63}{#1}}
\newcommand{\SpecialStringTok}[1]{\textcolor[rgb]{0.73,0.40,0.53}{#1}}
\newcommand{\StringTok}[1]{\textcolor[rgb]{0.25,0.44,0.63}{#1}}
\newcommand{\VariableTok}[1]{\textcolor[rgb]{0.10,0.09,0.49}{#1}}
\newcommand{\VerbatimStringTok}[1]{\textcolor[rgb]{0.25,0.44,0.63}{#1}}
\newcommand{\WarningTok}[1]{\textcolor[rgb]{0.38,0.63,0.69}{\textbf{\textit{#1}}}}
\usepackage{longtable,booktabs,array}
\newcounter{none} % for unnumbered tables
\usepackage{calc} % for calculating minipage widths
% Correct order of tables after \paragraph or \subparagraph
\usepackage{etoolbox}
\makeatletter
\patchcmd\longtable{\par}{\if@noskipsec\mbox{}\fi\par}{}{}
\makeatother
% Allow footnotes in longtable head/foot
\IfFileExists{footnotehyper.sty}{\usepackage{footnotehyper}}{\usepackage{footnote}}
\makesavenoteenv{longtable}
\setlength{\emergencystretch}{3em} % prevent overfull lines
\providecommand{\tightlist}{%
  \setlength{\itemsep}{0pt}\setlength{\parskip}{0pt}}
\usepackage{hyperxmp}
\usepackage{bookmark}
\IfFileExists{xurl.sty}{\usepackage{xurl}}{} % add URL line breaks if available
\urlstyle{same}
\hypersetup{
  pdftitle={The Journal as Substrate: Unifying Deployment{,} Replication{,} Backup{,} and Offline Operation in Distributed Systems},
  pdfauthor={Alvaro Rivera},
  pdfkeywords={\xmpquote{design theory}, \xmpquote{event
sourcing}, \xmpquote{journal-based
program}, \xmpquote{substrate}, \xmpquote{zero-downtime
deployment}, \xmpquote{cross-datacenter
replication}, \xmpquote{backup}, \xmpquote{offline
operation}, \xmpquote{replay}, \xmpquote{actor
systems}, \xmpquote{puppeteer framework}},
  hidelinks,
  pdfcreator={LaTeX via pandoc}}

\title{The Journal as Substrate: Unifying Deployment, Replication,
Backup, and Offline Operation in Distributed Systems}
\author{Alvaro Rivera}
\date{2026-05-14}

\begin{document}
\maketitle
\begin{abstract}
This paper is a design theory contribution in the sense of Alan Hevner,
Stuart March, Jinsoo Park, and Sudha Ram (2004). It identifies a
fragmentation that the canonical distributed-systems literature has
documented across four separate operational disciplines without
recognising them as instances of a single structural condition. It then
derives the property under which this fragmentation dissolves and
presents an instantiation demonstrating that the alternative is
realisable in production.

For decades, deployment without downtime, cross-datacenter replication,
backup, and offline operation have evolved as four parallel bodies of
practice, each with its own vocabulary, apparatus, and cost model. This
paper shows that this separation is contingent rather than structural.
Under the substrate condition established in prior papers of this series
--- where the journal is homoiconic, dense, and composed of compact
named operations rather than serialised states --- the four disciplines
reduce to a single primitive.

Deployment becomes replay against the substrate up to the present head.
Replication is the same replay at another site. Backup is the corpus
held without replay. Offline operation is replay after a delay.

The canonical apparatuses developed for these disciplines --- blue-green
environments, change-data-capture pipelines, snapshot agents, and
message-queue buffers --- perform work that the substrate already
provides by construction. Each arose independently to compensate for the
absence of the substrate condition rather than to extend it.

The construct introduced here names the substrate, formalises the four
equivalences that follow from it, and characterises the operational
regimes in which these equivalences can be observed.
\end{abstract}

\section{Operations from the
substrate}\label{operations-from-the-substrate}

\subsection{TL;DR}\label{tldr}

\begin{quote}
The canonical distributed-systems literature treats four operations ---
deployment without downtime, cross-datacenter replication, backup, and
offline catch-up --- as separate disciplines. Each has developed its own
mature apparatus: blue-green and rolling updates for deployment;
change-data-capture and log shipping for replication; snapshots and
point-in-time recovery for backup; message queues and convergence
protocols for offline operation. Systems that require all four must
implement and operate all four independently.

This paper shows that the separation is contingent. Under the
journal-as-substrate condition established in prior papers of this
series --- where the journal records operations rather than states, is
written in the same language as the program, and consists of compact
named entries --- the four disciplines reduce to a single primitive.

Deployment is replay to the present head. Replication is the same replay
at another site. Backup is the corpus held without replay. Offline
operation is replay after a delay. The canonical apparatuses become
structurally subsumed: each reconstructs, from outside the program, work
the substrate already performs by construction.

Six laboratories measure the substrate under conditions where these
equivalences hold. A version handover reaches the present head at 451k
entries/s; cross-site transfer carries 36--99.7 bytes per event against
payloads 3--9× larger; passive consumers converge across heterogeneous
backends (FileSystem, MySQL, SQL Server) to identical in-memory state in
18/18 cells; local append latency runs ≈3× below co-located RDBMS
engines under matched durability. §7 shows how four runtime primitives
--- each present for independent structural reasons --- compose into the
operations the canonical literature implements as separate subsystems.

\textbf{For a journaled program, deployment is replay, replication is
sharing history, backup is copying the program, and offline operation is
delayed replay.}
\end{quote}

\begin{center}\rule{0.5\linewidth}{0.5pt}\end{center}

\subsection{1. Introduction}\label{introduction}

This paper makes a design theory contribution. It identifies a
structural property whose absence has been implicitly assumed by
canonical distributed-systems literature, but never named as a single
construct; derives the equivalences that follow when the property is
present; and presents an instantiation showing it is realisable in
production. The contribution is conceptual; the instantiation serves as
an existence proof rather than as the substance of the claim. The genre
is that described by Alan Hevner, Stuart March, Jinsoo Park, and Sudha
Ram (2004) as design science research: an artefact satisfying the
construct's conditions, augmented by measurements that exhibit the
régime under which the construct holds.

This is neither a systems paper proposing a new mechanism nor a survey
of existing apparatuses. It does not introduce a deployment tool, a
replication protocol, a backup format, or an offline-operation pattern.
It examines a structural property that prior literature has implicitly
assumed absent and shows that, under its presence, four operational
disciplines of distributed systems reduce to a single primitive. In this
genre, contribution is measured by the precision of the construct, the
validity of the equivalences derived from it, and the realisability of
the instantiation. Sections 5 and 7 exhibit the instantiation and report
measurements establishing the régime in which the equivalences are
observable.

Distributed systems must address four operational dimensions: release
without downtime, cross-datacenter replication, recovery after data
loss, and periods of disconnectivity (offline operation). The canonical
approach treats these as separate disciplines, each with distinct
apparatus and cost semantics: blue-green and rolling updates;
change-data-capture and log shipping; snapshots and point-in-time
recovery; message-queue buffers and convergence protocols. These bodies
of practice are mature and widely adopted, yet they evolved largely
independently. Systems that require all four must implement and operate
four distinct stacks, each with its own engineering and cost.

The fragmentation is rarely questioned because it is familiar. But the
four disciplines can be restated more simply: they vary when and where a
program is read --- at the present head and in place (deployment), at
the present head at another site (replication), held at another site
without reading (backup), and read after a delay (offline catch-up).
They differ along two axes --- temporal position and spatial position
--- over a single artefact. If that artefact were the program itself,
recorded as operations rather than states, the four disciplines would be
four readings of the same primitive: reading the program in order.

This paper names that artefact and the equivalences it admits. The
construct introduced is the journal as the substrate of the program (in
the sense established in \href{02-program-value-separability.md}{Paper
2} §1.2: the pair of domain library and journal of invocations). Under
the anti-porosity condition of \href{01-anti-porosity.md}{Paper 1} and
the separability condition of
\href{02-program-value-separability.md}{Paper 2}, the journal records
operations rather than states, is written in the same language as the
program, and reduces each entry to a compact reference to a parametric
verb and its arguments. The journal, so constituted, is not a record
kept by the program; it is the program written out over time.

Four equivalences follow. Deployment is replay to the present head.
Replication is the same replay at another site. Backup is the corpus
held without replay. Offline operation is replay after a delay. The
canonical apparatuses above are responses to the absence of the
substrate condition; under its presence they become structurally
subsumed by the substrate.

§2 traces the genealogy of the four disciplines as parallel bodies of
literature. §3 names the assumption that produced the fragmentation. §4
states the substrate theorem and the four equivalences. §5 develops each
equivalence and reports measurements exhibiting the régime under which
it holds. §6 examines four canonical apparatuses --- blue-green
deployment, change-data-capture, snapshot backup, and message-queue
buffer --- and shows that each reconstructs, from outside the program,
work the substrate already performs. §7 presents an instantiation in
production. §8 relates the construct to prior papers in this series. §9
addresses counter-arguments. §10 concludes.

\begin{center}\rule{0.5\linewidth}{0.5pt}\end{center}

\subsection{2. Genealogy: four separate operational
disciplines}\label{genealogy-four-separate-operational-disciplines}

For decades, four operational concerns --- deployment downtime,
cross-datacenter replication, backup, and offline operation --- have
evolved as separate bodies of literature and practice. Each developed
its own vocabulary, its own apparatus, and its own cost model. The
genealogies traced below are not exhaustive; they are sufficient to
establish that the canon treats these four as distinct disciplines, each
with a dedicated mechanism, and that the resulting fragmentation of the
operational stack has not been challenged within any of the four
traditions.

\subsubsection{2.1 Deployment downtime}\label{deployment-downtime}

The earliest deployment practice in production systems was the cutover:
stop the running service, swap binaries, start it again. Downtime was
accepted as the cost of release. The first patterned response was
articulated by Fowler in \emph{BlueGreenDeployment} {[}Fowler 2010{]},
which named a procedure in which ``you ensure that you have two
production environments, as identical as possible'' and traffic is
switched from one to the other at release time. The mechanism is
environmental: two parallel deployments, one apparatus for routing, one
for binary distribution.

The continuous-delivery literature ratified the pattern. Humble and
Farley {[}Humble \& Farley 2010{]} generalized blue-green into a broader
release-engineering discipline in which deployment is a build-pipeline
output rather than an operator action, and downtime is a regression to
be detected rather than a default to be accepted.

The modern instantiations are canary release {[}Sato 2014{]} and the
orchestration-mediated rolling update {[}Burns et al.~2016; Kubernetes
documentation{]}. Canary release introduces partial traffic exposure
before full cutover, treating the deployed binary as a hypothesis tested
against a fraction of production load. The Kubernetes rolling update
replaces pods incrementally under a readiness-probe gate, making the
deployment apparatus a property of the orchestrator rather than of the
application.

A parallel strand of the deployment canon addresses schema and data
migration. Ambler and Sadalage {[}Ambler \& Sadalage 2006{]} formalized
\emph{database refactoring} as a discipline in which every change to the
data store is recorded as a forward-and-backward pair of scripts
versioned alongside the application code. The pattern is instantiated by
migration frameworks such as Rails Migrations, Liquibase {[}Liquibase
project{]}, and Flyway {[}Flyway project{]}, which apply versioned
scripts to the data store at release time. The migration apparatus is
separate from the traffic-switching apparatus: it runs adjacent to ---
or before --- the cutover, mutates the store in place, and produces a
schema or data shape that the next binary version expects.

Across these mechanisms the shape is constant: a separate apparatus is
built whose function is to absorb the change --- of binary, of schema,
or both --- without interrupting the service. The apparatus is external
to the program; the program is unaware that it is being replaced or
migrated.

\subsubsection{2.2 Cross-datacenter
replication}\label{cross-datacenter-replication}

Replication across datacenters originates in database engines. MySQL
replication {[}MySQL Reference Manual{]} established the canonical
primary/replica topology in the 1990s: the primary writes a binary log
of statements or row-level events, and replicas tail that log. The
replication mechanism lives inside the database engine; the application
program is unaware that replication is occurring.

The pragmatic maturation came with change-data-capture (CDC). Debezium
{[}Debezium project{]} and Maxwell {[}Maxwell project{]} read the
database's binary log and re-emit each change as an event on an external
transport, typically Kafka. CDC inserts an apparatus between the
database and downstream consumers: a process that watches state changes
in one system and reconstructs them as event streams for others. The
reconstruction step is necessary precisely because the originating
system does not emit events --- it emits state changes that CDC observes
and translates.

The modern instantiations extend the apparatus topologically.
Cassandra's multi-datacenter replication {[}Lakshman \& Malik 2010{]}
embeds the replication topology in the storage engine itself, allowing
writes accepted in any datacenter to propagate to peers under a
configurable consistency level. Kafka MirrorMaker {[}Kreps, Narkhede, \&
Rao 2011{]} generalizes the pattern for log-structured systems: a
dedicated process consumes from one cluster and re-produces to another,
with offsets tracking progress per topic.

Across these mechanisms the shape is constant: cross-datacenter
replication requires a separate apparatus --- engine-internal,
engine-adjacent, or fully external --- whose function is to observe the
changes happening in one location and reproduce them elsewhere.

\subsubsection{2.3 Backup}\label{backup}

The backup canon predates the others. Agent-based backup {[}Bacula
project; Veeam documentation{]} developed in the 1990s and 2000s as a
discipline in which a dedicated process periodically reads the
production data store and writes a copy to an archival medium. The agent
is external to the application; backup occurs at the storage layer.

The pragmatic maturation introduced snapshot semantics. The WAFL
filesystem {[}Hitz, Lau, \& Malcolm 1994{]} established the modern
primitive: a point-in-time copy taken at the filesystem level,
exploiting copy-on-write so that the snapshot is cheap and consistent
with respect to the moment of capture. Database engines adopted the same
pattern under the name \emph{point-in-time recovery} {[}PostgreSQL
documentation{]}, in which continuous archival of the write-ahead log
permits reconstruction of any committed state within the archived range.

The modern instantiations move the apparatus to object storage.
Cloud-era backup {[}AWS S3 Lifecycle; Azure Blob Storage backup
documentation{]} treats archival as a property of the storage tier
rather than of a dedicated agent: lifecycle policies move data to colder
tiers automatically, and recovery is a matter of retrieval from the
archive.

Across these mechanisms the shape is constant: backup is a separate
apparatus that produces a copy of state and is exercised at recovery
time to reconstruct that state. The backup procedure observes the data
store from the outside; the program being backed up does not
participate.

\subsubsection{2.4 Offline operation}\label{offline-operation}

The offline-operation canon is the most heterogeneous of the four,
because it spans network-layer, transport-layer, and application-layer
treatments of disconnection. The early precedent is store-and-forward
messaging: UUCP and, later, SMTP queue undeliverable messages locally
and retry until the remote endpoint is reachable. Disconnection is
absorbed by a queue at the boundary of the system.

The pragmatic maturation generalized the boundary queue into a
programmable apparatus. RabbitMQ {[}RabbitMQ documentation{]} introduced
dead-letter exchanges to capture messages whose delivery is impeded,
deferring them for retry or human inspection. Kafka {[}Kreps, Narkhede,
\& Rao 2011{]} generalized the pattern further: consumer offsets are
explicit, durable, and per-partition, so a consumer that goes offline
resumes from its last committed offset when it returns. The queue is no
longer a remediation mechanism for failed delivery; it is the substrate
over which the consumer's progress is tracked.

The modern instantiations carry the pattern into application-level
state. Eventually-consistent stores {[}DeCandia et al.~2007{]} treat
each replica as a local source of truth that converges with peers when
connectivity permits. The Dynamo paper established the design vocabulary
--- vector clocks, hinted handoff, read-repair --- that subsequent
systems (Cassandra, Riak, Redis replication {[}Redis documentation{]})
instantiate in varying degrees.

Once again, the shape is constant: offline operation requires a separate
apparatus --- a queue, an offset, a convergence protocol --- whose
function is to absorb the period during which the system cannot reach
its counterpart.

\subsubsection{2.5 Where the operational complexity
lives}\label{where-the-operational-complexity-lives}

The four disciplines surveyed above can be compared along a single
matrix. Each row names a discipline; each column names a dimension of
the apparatus the canon developed to handle that discipline.

{\def\LTcaptype{none} % do not increment counter
\begin{longtable}[]{@{}
  >{\raggedright\arraybackslash}p{(\linewidth - 6\tabcolsep) * \real{0.2500}}
  >{\raggedright\arraybackslash}p{(\linewidth - 6\tabcolsep) * \real{0.2500}}
  >{\raggedright\arraybackslash}p{(\linewidth - 6\tabcolsep) * \real{0.2500}}
  >{\raggedright\arraybackslash}p{(\linewidth - 6\tabcolsep) * \real{0.2500}}@{}}
\toprule\noalign{}
\begin{minipage}[b]{\linewidth}\raggedright
Discipline
\end{minipage} & \begin{minipage}[b]{\linewidth}\raggedright
Where the complexity lives
\end{minipage} & \begin{minipage}[b]{\linewidth}\raggedright
Apparatus
\end{minipage} & \begin{minipage}[b]{\linewidth}\raggedright
Cost
\end{minipage} \\
\midrule\noalign{}
\endhead
\bottomrule\noalign{}
\endlastfoot
Deployment & Outside the program, in the release pipeline &
Two-environment switch (blue-green); incremental pod replacement under a
readiness gate (rolling update); fractional traffic exposure (canary);
versioned schema and data migration scripts (Liquibase, Flyway, Rails
Migrations) & Duplicate production environment; orchestration layer;
manual or automated cutover procedure; forward-and-backward migration
pairs \\
Cross-datacenter replication & Inside the storage engine, or in an
adjacent CDC layer & Binary-log shipping (engine-internal);
change-data-capture (engine-adjacent); multi-DC topology (engine-aware)
& Replication lag; reconstruction overhead in CDC;
consistency-versus-availability configuration surface \\
Backup & Outside the program, in a dedicated archival pipeline &
Agent-based copy; filesystem or engine snapshots; continuous log
archival (PITR); object-store lifecycle & Recovery-time objective and
recovery-point objective; storage-tier cost; periodic exercise of the
recovery procedure \\
Offline operation & At the boundary between system and counterpart &
Persistent queue with retry semantics; durable consumer offsets;
convergence protocol over replicated stores & Buffer sizing; offset
management; eventual-consistency reasoning at the application layer \\
\end{longtable}
}

Three observations follow from the table. The location of the complexity
differs across rows: the deployment apparatus lives in the release
pipeline, the replication apparatus lives inside or beside the engine,
the backup apparatus lives in a separate archival pipeline, and the
offline-operation apparatus lives at the system's boundary. The
mechanism in each row was developed without reference to the others: the
blue-green literature does not cite the CDC literature; the snapshot
literature does not cite the consumer-offset literature. The four
literatures evolved in parallel, each solving what it believed to be a
different problem. And the cost in each row is paid separately: a system
that wishes to be zero-downtime, cross-datacenter, recoverable, and
tolerant of disconnection budgets independently for each of the four
columns.

The canon, taken as a whole, presents these four as four different
problems requiring four different apparatuses. §3 names the assumption
that made this fragmentation appear necessary.

\begin{center}\rule{0.5\linewidth}{0.5pt}\end{center}

\subsection{3. The assumption named}\label{the-assumption-named}

What §2 shows to be continuous across four separate bodies of literature
can be stated as a single structural assumption:

\begin{quote}
\textbf{Operational concerns of a system are addressed by apparatus
around the program rather than by the program itself.}
\end{quote}

Across the four disciplines, the response to each operational concern is
the introduction of a mechanism external to the program. Blue-green
deployment builds a parallel environment to absorb the cut-over.
Change-data-capture inserts a translation layer between stores to carry
replication. Snapshot agents observe the data store from outside and
write captures to archival media. Message-queue brokers hold messages at
the system boundary to absorb periods of disconnection. In each case,
the mechanism is not a sentence of the program it serves. It observes
the program from outside, maintains a parallel artefact, and is operated
independently of the program's own execution.

The assumption embedded in these practices is that operational concerns
require such external apparatuses, each developed independently to
address a specific condition.

The alternative this assumption leaves unexplored is that the program
might address these concerns by varying how it is read --- by changing
when or where the program is replayed, without leaving the program's own
representation. Under this alternative, the four disciplines become four
readings of a single primitive. The question becomes structural: under
what condition would the program admit such readings?

§4 names that condition and the four readings it permits. The
apparatuses surveyed in §2 remain appropriate responses in their
original context; what this paper shows is that they arise from the
absence of the substrate condition, which, if present, subsumes the need
for each apparatus.

\begin{center}\rule{0.5\linewidth}{0.5pt}\end{center}

\subsection{4. The substrate theorem}\label{the-substrate-theorem}

\subsubsection{4.1 The substrate}\label{the-substrate}

The first two papers in this series established two structural
properties of the journal that this paper takes as preconditions.
\href{01-anti-porosity.md}{Paper 1}, §3 named \emph{anti-porosity}: the
journal records the operations of the program rather than the states
they produce, and admits no representational sparsity at the boundary
between domain code and persisted form. The journal is therefore
homoiconic --- entries are sentences in the same language the program is
written in --- and dense --- every effect that crosses the boundary is
named in that language, not summarised as state delta.
\href{02-program-value-separability.md}{Paper 2}, §1.2 and §3 named
\emph{separability}: programs that are parametric in their arguments
admit a compiled form in which the operation's body is cached under an
identifier and the journal entry reduces to that identifier plus the
call's arguments. Together the two conditions yield a journal whose
entries are compact, named operations rather than serialised states.

Under these two conditions the journal is not a record kept by the
program; it is the program written out over time. Each entry is a
sentence the program has uttered; the journal is the corpus of those
sentences in the order they were uttered. The naming move of this paper
is to call this artifact \emph{the substrate of the program}: the
journal does not store the program's outputs, it constitutes the
program's existence over time. The program and the journal are the same
artifact under two readings --- one synchronous, one diachronic. Under
the synchronous reading, the runtime is the program. Under the
diachronic reading, the journal is the program.

\subsubsection{4.2 The theorem statement}\label{the-theorem-statement}

The theorem of this paper can be stated as a single proposition.

\begin{quote}
\emph{Given a program whose execution is recorded as a journal
satisfying the anti-porosity condition of Paper 1 and the separability
condition of Paper 2 --- that is, a journal that is homoiconic, dense,
and made of compact named operations rather than serialised states ---
the following four equivalences hold.}
\end{quote}

\textbf{E1 --- Deployment is replay.} A new version of the runtime,
started from the substrate, reconstructs the program by reading the
journal sequentially and applying each entry. Nothing about the new
version's instantiation differs in kind from what the previous version
did at every entry up to that point; the act of bringing a process to a
state in which it can serve requests is the act of replaying the
substrate up to its present head. Deployment therefore differs from
steady-state operation only in the position of the read cursor over the
same corpus.

\textbf{E2 --- Replication is sharing history.} A second instance of the
program, running on a different machine or in a different site,
reconstructs the program by reading the same sequence of entries the
first instance wrote. The act of replicating is the act of transmitting
the substrate from one site to another and replaying it at the
destination; no state extraction, no schema mapping, and no change-data
inference enter the act. Replication therefore differs from deployment
only in the site at which the same corpus is replayed.

\textbf{E3 --- Backup is copying the program.} A consumer that reads
entries from the substrate and persists them locally --- without
instantiating the runtime that would replay them --- holds a copy of the
program in the form in which it was written. The act of backing up is
the act of duplicating the substrate; recovery is the act of replaying
the duplicate. Backup therefore differs from replication only in whether
the destination chooses to replay the same corpus or merely to hold it.

\textbf{E4 --- Offline operation is delayed replay.} A consumer that is
unreachable when entries are written, and that catches up by reading
them later from a persistent buffer or directly from the source, applies
the same replay to the same sequence as a consumer that received the
entries in real time. The act of operating offline is the act of
replaying the substrate after a delay; the delay does not enter the
program, only the cursor of the consumer. Offline operation therefore
differs from replication only in the elapsed time between the write and
the read of the same corpus.

The four equivalences are not parallel restatements of a single
mechanism. Each addresses one classical operation of distributed systems
(§2) and shows that the operation, as the canon presents it, dissolves
into a single primitive --- replay --- once the substrate's properties
are taken as given. E1 addresses the time axis (a new version replacing
an old one in place); E2 addresses the space axis (a new site running
the same program); E3 addresses the storage axis (a consumer holding the
corpus without running it); E4 addresses the connectivity axis (a
consumer reading later than the writer wrote). The four together exhaust
the operational variations that the canonical literature in §2 treats as
separate disciplines. The operational specifics of how each replay is
realised --- handoff, transmission, passive ingestion, catch-up --- are
taken up in §5; §4 confines itself to the structural claim.

\subsubsection{4.3 Apparatus: the substrate and its four replay
arrows}\label{apparatus-the-substrate-and-its-four-replay-arrows}

The shape of the theorem can be rendered as a simple spatial
arrangement. No new content is introduced; the arrangement makes the
four equivalences of §4.2 visible at a glance, with the substrate at the
centre and the four conditions arrayed around it.

The substrate --- the journal as program --- sits at the centre. From it
radiate four equivalences, each corresponding to one classical
operational discipline:

\begin{itemize}
\tightlist
\item
  \textbf{E1 --- Deployment}: replay to the present head.
\item
  \textbf{E2 --- Replication}: replay of the same corpus at another
  site.
\item
  \textbf{E3 --- Backup}: the corpus held without replay.
\item
  \textbf{E4 --- Offline operation}: replay after a delay.
\end{itemize}

The vertical pair (E1, E4) varies the temporal position of the replay
--- at the present head, or after a delay. The horizontal pair (E2, E3)
varies the spatial position of the corpus --- replayed at a different
site, or merely held at one. The canonical literature in §2 names four
disciplines by looking at the destinations; viewed from the substrate,
only replay occurs, under four conditions on when and where.

\begin{center}\rule{0.5\linewidth}{0.5pt}\end{center}

\subsection{5. Consequences: four equivalences as variants of
replay}\label{consequences-four-equivalences-as-variants-of-replay}

Each of the four sub-sections that follow develops one equivalence of
§4.2. The sub-section opens with the operational realisation of the
equivalence --- the structural intervals through which the program is
read under that variant --- and follows with a measurement that exhibits
the régime in which the equivalence holds. Each lab is designed to
isolate the structural content of its equivalence, not to optimise
throughput; performance figures are reported as evidence of the régime,
not as benchmark claims.

\subsubsection{5.1 E1 --- Deployment is
replay}\label{e1-deployment-is-replay}

§4.2 stated the equivalence: a new version of the runtime, started from
the substrate, reconstructs the program by reading the journal
sequentially and applying each entry. §5.1 develops the operational
realization of this equivalence --- the handover window during which an
instance about to become authoritative reads the substrate up to the
present head, and the precise sense in which that window is replay and
nothing else.

\paragraph{The handover window}\label{the-handover-window}

A new version of the runtime begins as a \emph{follower}: a process that
reads from the same substrate the live instance writes but does not yet
accept requests. The follower's task is to reach the present head of the
substrate. Once it has --- and once a brief synchronization tail has
cleared --- the follower is admitted as authoritative and the previous
version is released. Three structural intervals comprise the window:

\begin{enumerate}
\def\labelenumi{\arabic{enumi}.}
\tightlist
\item
  \textbf{Bulk replay.} The follower opens the journal from entry zero
  (or from a checkpoint, if the substrate is partitioned with one) and
  applies each entry to its in-memory state in the order written. The
  state reconstructed at the end of bulk replay equals, entry-for-entry,
  the state the previous version reached by writing those same entries.
\item
  \textbf{Lock and synchronization tail.} The previous version is paused
  briefly so that no further entries are written. The follower reads the
  small residue of entries that arrived between the start of bulk replay
  and the pause --- the \emph{handover tail} --- and applies them. The
  state at the end of the tail equals the state the previous version
  held at the moment it was paused.
\item
  \textbf{Gate flip.} The follower is admitted as the authoritative
  instance. From this point on the substrate's writer is the new
  version; the previous version is released.
\end{enumerate}

The operational specifics of each interval are realized by a single
primitive in the instantiation (§7.1 develops the gate mechanism); what
§5.1 establishes is that the \emph{content} of the window is replay. No
state extraction, no schema migration, and no inference step enters the
operation. The new version becomes authoritative by reading what the
previous version wrote, in the order it was written, until the read
cursor reaches the position the previous version had left it at.

\paragraph{What the substrate
measures}\label{what-the-substrate-measures}

L1 (lab
\href{data/paper05-lab1-redblack-replay-time/}{\texttt{paper05-lab1-redblack-replay-time}}
@ \texttt{ef3a002}) instruments the handover window over four journal
sizes (N ∈ \{1k, 10k, 100k, 1M\} entries) under the compact-action
régime named in Paper 2. The measurement isolates the inner bulk-replay
interval and the handover tail separately so that the structural content
of the window --- the act of replay --- is reported in its own right
rather than aggregated with orchestration overhead.

Under the compact-action régime with \texttt{AlwaysCompiled}
compilation, the substrate replays N entries at a sustained rate of
\textbf{451 thousand entries per second} (p50, N = 100 000), completing
the bulk replay in \textbf{221.6 ms} (p95 = 237.4 ms). The rate is
reached as the runtime amortises between N = 10 000 and N = 100 000
entries and holds through the N = 1 000 000 anchor at 416 thousand
entries per second. The handover tail is below 30 µs at every N
measured. \textbf{Bulk replay accounts for more than 99.8 \% of the
deploy window in every cell.} The replay rate is sustained across a
tenfold journal range; were it linear in expanded state rather than in
compact-event count, the rate would fall as state grew. It does not.

Three properties of the measurement deserve naming. First, the substrate
cost of deployment is linear in compact-event count, not in expanded
state --- replaying ten times the journal takes ten times the time, at
constant rate. Second, the orchestration tail is empirically negligible
at every measured size; the deploy window \emph{is} the replay window to
within a fraction of one percent. Third, the régime under which this
holds is the régime Paper 2 already named the canonical one: parametric
verbs encoded as \texttt{ActionId} plus arguments, not literal scripts.
A journal of literal scripts would replay slower in proportion to script
length; this paper concerns the compact-action substrate named in Paper
2.

\paragraph{Apparatus}\label{apparatus}

The handover window can be rendered as a sequence between four roles.
The leader is the live instance accepting writes; the journal is the
substrate; the follower is the new instance approaching the present
head; the gate is the predicate that decides which instance is
authoritative at any moment.

\begin{Shaded}
\begin{Highlighting}[]
\NormalTok{sequenceDiagram}
\NormalTok{    autonumber}
\NormalTok{    participant L as Leader (authoritative)}
\NormalTok{    participant J as Journal (substrate)}
\NormalTok{    participant F as Follower (new version)}
\NormalTok{    participant G as Gate}

\NormalTok{    Note over F,G: Follower starts; gate closed for follower}
\NormalTok{    F{-}\textgreater{}\textgreater{}J: open from entry 0 (or checkpoint)}
\NormalTok{    loop bulk replay}
\NormalTok{        J{-}{-}\textgreater{}\textgreater{}F: next entry}
\NormalTok{        F{-}\textgreater{}\textgreater{}F: apply entry to in{-}memory state}
\NormalTok{    end}
\NormalTok{    Note over F: follower at near{-}head}
\NormalTok{    F{-}\textgreater{}\textgreater{}L: request pause}
\NormalTok{    L{-}\textgreater{}\textgreater{}L: stop accepting writes}
\NormalTok{    loop handover tail}
\NormalTok{        J{-}{-}\textgreater{}\textgreater{}F: residual entry written before pause}
\NormalTok{        F{-}\textgreater{}\textgreater{}F: apply}
\NormalTok{    end}
\NormalTok{    Note over F: follower at present head}
\NormalTok{    F{-}\textgreater{}\textgreater{}G: open for follower}
\NormalTok{    G{-}{-}\textgreater{}\textgreater{}L: close for leader}
\NormalTok{    Note over L: leader released; follower authoritative}
\end{Highlighting}
\end{Shaded}

The diagram contains no operation other than reading from the journal
and applying. The leader writes nothing between pause and release; the
follower reads from the same artifact at the same offsets the leader was
about to read had it continued. The gate's flip is bookkeeping: it
nominates which instance is authoritative; it neither produces nor
consumes substrate content.

\paragraph{One canonical exception}\label{one-canonical-exception}

The equivalence developed here covers the in-process deployment of a
single program. One canonical case lies outside its scope: a coordinated
cut-over of a producer/consumer pair on an external message broker
(typically Kafka) when the message format itself changes between
versions. There the two endpoints must release together, because the
substrate they share is the broker's topic --- a contract between
independent programs --- rather than the journal of either. The
construct of this paper does not deny this case; it narrows its scope.
Within a journaled program, deployment is replay; across two journaled
programs that share a broker topic with a changing schema, the cut-over
remains a separate concern. Paper 3, claim 8, named this exception in
the discussion of in-actor continuity; it is preserved here without
amendment.

\paragraph{To the next equivalence}\label{to-the-next-equivalence}

§5.1 has shown that bringing a new version of a program to a state in
which it can serve is replay of the substrate up to the present head,
and that the orchestration around that replay is residual at every
measured size. The equivalence varies the position of the read cursor in
time --- at the present head of the same substrate. §5.2 varies the
position of the substrate in space: a second instance, at another site,
reconstructs the same program by reading the same sequence of entries
the first instance wrote.

\subsubsection{5.2 E2 --- Replication is sharing
history}\label{e2-replication-is-sharing-history}

§4.2 stated the equivalence: a second instance of the program, running
on a different machine or in a different site, reconstructs the program
by reading the same sequence of entries the first instance wrote. §5.2
develops the operational realization in two parts. First, that the bytes
streamed between sites are structurally dense --- the substrate's wire
form carries named operations, not serialized states, by construction.
Second, that a consumer reading those bytes reaches the same program as
the writer by replay alone, without recourse to a coordinator. The
transport across which the bytes travel is treated as a witness --- an
off-the-shelf service that satisfies the substrate's requirement on
consumer-pull delivery.

\paragraph{The compact wire}\label{the-compact-wire}

\href{02-program-value-separability.md}{Paper 2}, §2.3 (\emph{dense
journaling as reference}), established that under separability the
journal entry of an invocation reduces to a reference to the parametric
verb's definition plus the arguments of the call. The verb's body is not
present in the entry; it lives once, named, in the corpus's vocabulary.
Replication across datacenters is the transmission of those entries, not
of states; the wire encodes the program as it was uttered, not as it was
applied.

L2 (lab
\href{data/paper05-lab2-bytes-per-event/}{\texttt{paper05-lab2-bytes-per-event}}
@ \texttt{c68b2f4}) instruments the codec on the substrate's wire
surface, comparing the compact encoding (\texttt{ActionId} + parameters)
against the literal encoding (script DSL + parameter assignments) at
three tiers of verb richness, with and without an opportunistic
\texttt{gzip} on the literal form.

Bytes-per-event in compact form is 36 at the trivial-arithmetic tier
(68-byte script body), 36 at the branching-arithmetic tier (99-byte
body), and 99.7 at the production-shaped tier (681-byte body). The
literal form at the same three tiers is 121, 153, and 913.7 bytes. The
ratio between literal and compact rises from \textbf{3.4× at tier 1} to
\textbf{9.2× at tier 3}. Applying \texttt{gzip} to the literal form
closes the gap only to 3.0× / 3.7× / \textbf{3.9× at tier 3}. The
structural saving cannot be recovered by an opportunistic compressor,
because the script body is not present in the compact form at all --- it
lives once in the definition record (912 bytes at tier 3, amortized to
≈0.91 bytes per invocation over 1 000 calls) and is referenced by a
4-byte \texttt{ActionId} thereafter.

The cross-validation against Paper 2 Lab 4 sharpens the claim. Tier 3's
compact payload is 67.7 bytes here; Paper 2 Lab 4 measured 67 bytes for
the production verb that the synthetic tier-3 stand-in calibrates
against. The two numbers agree within rounding. At full production
scale, \href{02-program-value-separability.md}{Paper 2} Lab 4 reports
the literal-to-compact ratio reaches 20× over a thousand invocations of
a real ticket-purchase verb. The 9.2× headline here is a conservative
lower bound under the lab's literal-projection; the production figure is
reported in Paper 2 Lab 4.

Three properties of the measurement deserve naming. First, density is
structural rather than algorithmic --- compression operating after the
fact cannot recover what naming achieved by construction. Second, the
saving grows with verb richness; a domain whose operations are named
transitions (\href{02-program-value-separability.md}{Paper 2} §3)
replicates per event at a width that does not scale with the operation's
internal complexity. Third, replication across datacenters, in this
régime, is event streaming over a substrate that does not carry state by
construction --- the wire is the program as uttered, not as expanded.

\paragraph{The push-to-pull inversion}\label{the-push-to-pull-inversion}

The substrate is written by the producer and read by the consumer; the
role of the transport between them is to hold the corpus durably and to
deliver entries at the consumer's pace. The producer does not know the
consumer's rate; the consumer's progress is measured against the
substrate's own offsets, not against any signal from the producer. A
consumer that lags absorbs the lag in its own cursor; the producer's
write path is unaffected.

This inversion is not a property of any particular delivery service. It
is a property of the substrate. The producer's journal is the source of
record; any consumer downstream is a reader of that record. An
off-the-shelf webhook delivery service satisfies the consumer-pull
requirement without adding any structure the substrate did not already
license: it accepts an inbound POST from the producer's local journal
callback, holds it durably, and offers it to subscribers at the rate
they fetch. The fact that an off-the-shelf service suffices is the
relevant observation; the service is a witness to the substrate's
structural sufficiency, not an extension of it.

The shape can be rendered between five roles: the primary instance at
site A, the primary's journal, the delivery service, the consumer at
site B, and the consumer's journal.

\begin{Shaded}
\begin{Highlighting}[]
\NormalTok{sequenceDiagram}
\NormalTok{    autonumber}
\NormalTok{    participant P as Primary (site A)}
\NormalTok{    participant Ja as Journal A (substrate)}
\NormalTok{    participant T as Delivery (off{-}the{-}shelf, e.g. SVIX)}
\NormalTok{    participant C as Consumer (site B)}
\NormalTok{    participant Jb as Journal B (substrate)}

\NormalTok{    P{-}\textgreater{}\textgreater{}Ja: append entry n}
\NormalTok{    Ja{-}\textgreater{}\textgreater{}P: callback (entry n written)}
\NormalTok{    P{-}\textgreater{}\textgreater{}T: post entry n}
\NormalTok{    Note over T: held durably; offered at subscriber\textquotesingle{}s pace}
\NormalTok{    C{-}\textgreater{}\textgreater{}T: fetch next}
\NormalTok{    T{-}{-}\textgreater{}\textgreater{}C: entry n}
\NormalTok{    C{-}\textgreater{}\textgreater{}Jb: append entry n}
\NormalTok{    C{-}\textgreater{}\textgreater{}C: apply entry to in{-}memory state}
\end{Highlighting}
\end{Shaded}

The diagram contains no synchronous coupling between producer and
consumer. The producer's append-and-post is local; the consumer's
fetch-and-apply is local; the durable buffer between them is the
delivery service. The substrate at A and the substrate at B contain the
same entry at the same offset --- the two journals are bit-equal under
the measurement examined next.

\paragraph{The symmetric consumer}\label{the-symmetric-consumer}

§4.2 stated that replication ``differs from deployment only in the site
at which the same corpus is replayed.'' The structural content of this
claim is that the consumer at site B, running the same binary as the
producer at site A, derives the same program from the corpus it reads as
the producer derived from the corpus it wrote. No coordinator is
required to synchronise the two sites; the substrate is the coordinator.
This property removes the need for any external coordinator between
sites: coordination is implicit in the shared corpus rather than enacted
by a privileged runtime.

L3 (lab
\href{data/paper05-lab3-inproc-symmetric/}{\texttt{paper05-lab3-inproc-symmetric}}
@ \texttt{3984c5d}) measures this property at the level of Reactions ---
the partition primitive of \href{03-reactions-and-partition.md}{Paper 3}
--- under two terminators, \texttt{Emit} and \texttt{MarkAsSkip}
(Elide). Two instances of the same actor binary consume the same event
stream; the lab compares, at the byte level, both the callback tuples
each instance's Reactions produce and the journal segments each instance
writes. Across 6 cells (N ∈ \{100, 500, 1 000\} × K = 2 repetitions),
\textbf{all 6 cells produce 0 callback byte diffs and 0 journal-segment
byte diffs}. The instance at site B reconstructs the Reaction output
entirely from the journal prefix it reads.

The discipline that makes this work without a coordinator is a property
of the program, not of the runtime. A Reaction's guard refers to data
the program has explicitly \emph{exposed} alongside its events --- the
\texttt{expose} clause established in
\href{03-reactions-and-partition.md}{Paper 3}. Whatever a Reaction at B
needs to evaluate a guard travels in the corpus the program writes;
whatever the program chose not to expose, no remote evaluator can see.
The substrate carries the program, and the program carries the contract
on which its consumers depend.

\href{04-cross-actor-continuity.md}{Paper 4}, §5.3 and claim 10, named
the closely related property for cross-actor causation: the journal
record of the act of speaking from one actor to another is what survives
cross-datacenter replication, where canonical apparatuses (CDC, log
shipping, message-queue replication) reconstruct state and lose the
causal chain that produced it. §5.2 of the present paper completes the
picture from the substrate side: the journal record is what the consumer
reads, and the consumer's Reactions derive from it without needing the
producer's runtime to be reachable.

Two caveats are worth naming, in the spirit of Capa 2 honesty. L3's
transport is in-process --- a channel between two instances in the same
process --- and its measurement caps at N = 1 000 entries (beyond which
the harness becomes memory-bound). The symmetric property is
structurally insensitive to the transport's distance and to N: a
Reaction's match path is per-entry, and the journal-segment diff at N =
1 000 is identically zero across all measured cells. But the lab
demonstrates the property in vitro; full cross-datacenter delivery is
the operational mode realised in §7.2, and the higher-N regime is
bounded by the consumer's append throughput rather than by anything the
matcher does.

\paragraph{What §5.2 establishes, and what §5.3 develops
next}\label{what-5.2-establishes-and-what-5.3-develops-next}

§5.2 has shown that replication, in the substrate's terms, is the
streaming of named operations over a wire whose density is by
construction; that the transport between sites can be an off-the-shelf
consumer-pull service; and that a consumer running the same binary
reaches the same program as the producer by reading the corpus it
shares. The equivalence varies the spatial position of the corpus while
preserving replay as the means of reconstruction.

§5.3 varies the spatial position of the corpus without reconstructing
the program at the destination. A consumer that reads the substrate and
persists it locally --- without instantiating the runtime that would
replay it --- holds the program in the form it was written. Backup, in
this view, is replication arrested at the moment of holding rather than
carried through to replay.

\subsubsection{5.3 E3 --- Backup is copying the
program}\label{e3-backup-is-copying-the-program}

§4.2 stated the equivalence: a consumer that reads entries from the
substrate and persists them locally --- without instantiating the
runtime that would replay them --- holds a copy of the program in the
form in which it was written. The canonical literature on backup (§2.3)
treats backup as a captured state, archived periodically by an external
apparatus and replayed only on recovery. The equivalence stated here
re-reads the operation: for a journaled program, backup is the program
materializing in another substrate, and the destination's choice is
whether to hold the corpus or replay it.

\paragraph{The structural argument}\label{the-structural-argument}

The equivalence follows in four steps from the substrate condition of
§4.1:

\begin{enumerate}
\def\labelenumi{\arabic{enumi}.}
\tightlist
\item
  If the journal is the program (Papers 1 and 2), the artifact whose
  duplication preserves the program is the journal.
\item
  Duplication of the journal is the streaming of named entries to a
  destination substrate --- the program materializing there, sentence by
  sentence, in the order it was uttered.
\item
  The destination need not instantiate the runtime to receive the
  corpus; it need only persist what arrives. The consumer is a
  substrate, not a process.
\item
  When and if the destination chooses to replay, the corpus reconstructs
  the program. The choice between \emph{holding} and \emph{replaying} is
  independent of the act of materializing.
\end{enumerate}

\paragraph{Three symmetries}\label{three-symmetries}

Backup, in this view, is positioned by its relation to the three other
equivalences of §4.2. The position is fixed by what the destination
commits to, not by any new operation.

\begin{itemize}
\tightlist
\item
  \textbf{Backup and replication.} Replication (§5.2) is the destination
  streaming the corpus and replaying it; backup is the destination
  streaming the corpus and holding it. The same protocol governs both.
  Replication is backup carried through to replay; backup is replication
  arrested before replay.
\item
  \textbf{Backup and deployment.} Deployment (§5.1) is replay at a
  different time of the same site, ending with the new instance admitted
  as authoritative; backup is replay-eligible material held at a
  different site, with no instance ever admitted as authoritative there.
  The same protocol reads the corpus forward; what differs is the gate
  flip at the end --- present in deployment, absent in backup.
\item
  \textbf{Backup and offline operation.} Offline operation (§5.4) is the
  consumer reading the corpus with a delay between write and read but
  eventually catching up and serving; backup is the consumer reading the
  corpus and never serving. The same protocol delivers entries; what
  differs is whether the consumer ever takes a turn at serving requests.
\end{itemize}

The three symmetries are not corollaries discovered after the fact. They
follow from the substrate condition directly: each of the four
equivalences is read off the substrate by varying one of two axes ---
when the read happens, and what the consumer commits to once it has read
--- and the four exhaust the combinations. §4.3 named this arrangement.

\paragraph{What the substrate
measures}\label{what-the-substrate-measures-1}

L4 (lab
\href{data/paper05-lab4-passive-consumer/}{\texttt{paper05-lab4-passive-consumer}}
@ \texttt{a21a655}) instruments the materialize cycle on three storage
backends under the compact-action régime of L1. The dataset spans 3
backends × 3 journal sizes × 2 protocol layers × K = 2 measurement
repetitions = 36 cells, plus 3 catch-up cells.

Under régime N = 100 000 compact entries, the destination journal
converges at \textbf{≈ 2 419 events/sec on FileSystem}, \textbf{≈ 1 004
events/sec on MySQL}, and \textbf{≈ 764 events/sec on SQL Server (SQL
Edge stand-in)}. The wire round-trip itself accounts for \textbf{less
than 0.6 \% of the cycle} at every cell; the substrate cost is the
destination's append rate, not the operation. A replica actor
instantiated over the destination journal --- a verification step, not
part of the backup operation --- reaches the same in-memory state as the
primary in \textbf{18 of 18 measured cells across the three backends and
three journal sizes}.

The independence from backend is the empirically significant property
here. The equivalence is not a feature of any storage technology; it
holds wherever the journal can be appended to.

\paragraph{The destination as cursor}\label{the-destination-as-cursor}

The destination's three reading regimes --- near the head of the
substrate, paused after an interruption, or newly instantiated after the
primary has run for some time --- share a single protocol. Each is a
cursor on the substrate, read forward from its present position to the
present head. Catch-up after a retention loss is not a separate mode; it
is the same read, started further back. L4's catch-up cells confirm
this: the catch-up throughput matches the steady-state throughput within
stochastic variation on each backend. The substrate does not contain a
distinction between \emph{current} and \emph{recovering}; it contains a
sequence of entries, and each consumer holds a cursor over it.

\paragraph{Caveats}\label{caveats}

Two caveats are worth naming, in the spirit of Capa 2 honesty. First,
L4's transport between primary and destination is in-process --- a local
proxy, not a real cross-datacenter wrapper; the network latency × number
of round-trips a cross-datacenter deployment adds is bounded but not
measured here. Second, L4 measures the \emph{passive} half of the
equivalence --- the destination consuming the corpus --- not the act of
promoting a destination to authoritative status; that promotion is the
E1 case §5.1 already developed, and is structurally available to any
destination that holds the corpus to a chosen EntryId.

The operational realisation of E3 --- the modes under which a
destination starts, how the program declares which entries materialize
where, the transport that carries them, and the dimensions along which
the construct scales (heterogeneous backends, multi-destination fanout,
per-datacenter topology) --- is developed in §7.3.

\paragraph{To the next equivalence}\label{to-the-next-equivalence-1}

§5.3 has shown that backup is the program materializing in another
substrate; that backup is positioned by its relation to the three other
equivalences along two structural axes --- when the read happens, and
what the consumer commits to once it has read; and that the corpus
convergence rate on heterogeneous backends is independent of the
construct. The equivalence varies the consumer's commitments while
preserving replay as the means by which the program is reconstructed if
and when needed.

§5.4 varies a different axis: the delay between when the substrate is
written and when it is read. Offline operation, in this view, is replay
after a gap --- the consumer reads the same corpus the producer wrote,
but the read trails the write by an arbitrary interval.

\subsubsection{5.4 E4 --- Offline operation is delayed
replay}\label{e4-offline-operation-is-delayed-replay}

§4.2 stated the equivalence: a consumer that is unreachable when entries
are written, and that catches up by reading them later from a persistent
buffer or directly from the source, applies the same replay to the same
sequence as a consumer that received the entries in real time. The
canonical literature on offline operation (§2.4) treats the buffer as an
apparatus added at the system's boundary --- a message queue, a durable
offset, a convergence protocol --- whose function is to absorb the
period during which the system cannot reach its counterpart. In the
regime described here, no such apparatus exists: the substrate itself is
the durable buffer, and offline operation reduces mechanically to cursor
advancement over the same corpus.

If the producer writes to a substrate that is local --- local to the
producer's process, locally durable --- then there is no point at which
the producer needs to know whether any consumer is reachable. The
producer's write is to its own substrate. A consumer reading from that
substrate, whether it reads in real time or after an interruption, holds
a cursor over the same corpus; when the cursor advances, the consumer
applies the entries between the old position and the new one. The buffer
of the canonical literature is therefore not added; it is the substrate
itself.

\paragraph{Offline as the limit of
replication}\label{offline-as-the-limit-of-replication}

§5.2 developed replication as the consumer reading the corpus the
producer wrote, at a different site, with the transport between them
holding entries durably until the consumer fetches them. §5.4 is the
limit case: the consumer's read trails the write not by transport
latency but by an arbitrary interval --- minutes, hours, or longer. The
protocol is identical. The act of catching up is the act of advancing
the cursor over the entries that accumulated in the meantime.

The same mechanism applies in the reverse direction. A producer whose
canonical remote store is unreachable continues writing to its local
substrate; the remote store, when reachable again, is a consumer that
fell behind. The role of \emph{producer} and \emph{downstream store} is
asymmetric only in steady-state framing; structurally, both are cursors
over the same corpus.

\paragraph{What the substrate
measures}\label{what-the-substrate-measures-2}

L5 (lab \href{data/paper05-lab5-offline/}{\texttt{paper05-lab5-offline}}
@ \texttt{d5cb906}) measures the offline equivalence by configuring an
actor's substrate as a local persistent journal with an asynchronous
flush to a remote canonical backend (MySQL and Azure SQL Edge). The
remote backend is then partitioned for a measured interval
(\textasciitilde5 s of write volume); the partition is resolved; the
backlog drains.

Three properties of the measurement deserve naming.

\begin{enumerate}
\def\labelenumi{\arabic{enumi}.}
\tightlist
\item
  \textbf{Decoupling of write latency from the remote.} Under online
  operation, the primary appends at \textbf{517 µs p50 against MySQL}
  and \textbf{455 µs p50 against SQL Edge} --- \textbf{7.29× and 2.92×
  faster} than the unbuffered direct path. The producer's lock release
  time is bounded by the local substrate's append, not by the remote's
  commit latency.
\item
  \textbf{Partition latency is not worse than online latency.} During
  the partition, the buffered primary's per-append p50 is 338 µs on
  MySQL and 381 µs on SQL Edge --- the same régime as online; the write
  lock never traverses the network. The substrate's local persistence is
  what makes this hold; the absence or presence of the remote does not
  enter the write path.
\item
  \textbf{Catch-up is linear in the backlog.} On reconnect, the
  accumulated entries drain at the remote backend's steady-state ingest
  rate --- ≈ 280 events/sec on MySQL, ≈ 420 events/sec on SQL Edge ---
  with \textbf{zero events lost} across the measured cells. Catch-up is
  not a separate mode; it is the remote's cursor advancing from its last
  position to the present head.
\end{enumerate}

Two caveats are worth naming, in the spirit of Capa 2 honesty. First, L5
partitions the remote by stopping its container --- a realistic
partition mode but not the only one; half-open connections without
container death would exercise the connection-timeout path differently
and are not measured here. Second, the partition-phase latency being at
or below the online phase is partly an artifact of single-host
scheduling --- the asynchronous flush thread is idle during partition,
freeing CPU for the writer; in production with separate threads on
separate cores, the artifact shrinks, but the structural property ---
the write path is independent of the remote's reachability --- holds
either way.

The operational realisation of E4 --- the configuration that enables the
local substrate as the primary's source of truth, the asynchronous
component that flushes to the canonical backend, and the catch-up path
on reconnect --- is developed in §7 alongside the other operational
details.

\paragraph{To the next equivalence}\label{to-the-next-equivalence-2}

§5.4 has shown that offline operation is replay after a gap, and that
the gap is absorbed by the substrate's own local persistence rather than
by a separate apparatus. The four equivalences of §4.2 are now
developed.

§5.5 turns to the cost that underlies all four: replay against a local
substrate is cheap because append against a local substrate is cheap.
The latency budget that supports the four equivalences is examined as a
single property of the substrate, not as a property of any one
operation.

\subsubsection{5.5 The latency budget}\label{the-latency-budget}

The four equivalences of §4.2 reduce four classical operational
disciplines to a single primitive: replay. Replay is the act of reading
the substrate forward from a cursor and applying each entry. Its cost is
the cost of reading and applying entries. For a substrate that is read
often and written often, the cost of \emph{writing} entries ---
appending to the journal --- is the cost basis that supports the family
of replay-based operations. §5.5 examines that basis.

The structural claim that underlies E1--E4 is this: every replay-based
operation reduces to advancing a cursor over N entries, and therefore
inherits the substrate's per-entry append latency as its dominant cost
term. The actor's apply cost is identical across all regimes because it
operates on the same in-memory state. The only term that varies between
regimes is the cost of appending entries to the substrate.

The relevant comparison is between two regimes for the substrate's
persistence. Under the regime this paper develops, the substrate is a
local journal whose persistence is one \emph{fsync}-bounded append per
entry. Under the canonical regime (§2), the substrate's role is played
by an RDBMS whose persistence is one commit-bounded \emph{INSERT} per
entry. The latency gap between the two regimes propagates linearly to
the cost of deployment, replication, backup, and offline catch-up. The
four operational disciplines therefore share a single cost basis: the
substrate's append latency.

\paragraph{What the substrate
measures}\label{what-the-substrate-measures-3}

L6 (lab
\href{data/paper05-lab6-latency-budget/}{\texttt{paper05-lab6-latency-budget}}
@ \texttt{86905dc}) measures per-entry append latency on three local
backends under a one-event = one-durable-commit régime: a local
FileSystem journal, a co-located SQL Server (Azure SQL Edge stand-in),
and a co-located MySQL. The dataset spans 3 backends × K = 5 runs × N =
100 000 entries per backend = 1.5 million samples, with the first 1 000
appends per cell discarded as warm-up.

{\def\LTcaptype{none} % do not increment counter
\begin{longtable}[]{@{}
  >{\raggedright\arraybackslash}p{(\linewidth - 8\tabcolsep) * \real{0.2596}}
  >{\raggedright\arraybackslash}p{(\linewidth - 8\tabcolsep) * \real{0.1635}}
  >{\raggedright\arraybackslash}p{(\linewidth - 8\tabcolsep) * \real{0.1635}}
  >{\raggedright\arraybackslash}p{(\linewidth - 8\tabcolsep) * \real{0.1635}}
  >{\raggedright\arraybackslash}p{(\linewidth - 8\tabcolsep) * \real{0.2500}}@{}}
\toprule\noalign{}
\begin{minipage}[b]{\linewidth}\raggedright
Backend
\end{minipage} & \begin{minipage}[b]{\linewidth}\raggedright
Append p50 (µs)
\end{minipage} & \begin{minipage}[b]{\linewidth}\raggedright
Append p95 (µs)
\end{minipage} & \begin{minipage}[b]{\linewidth}\raggedright
Append p99 (µs)
\end{minipage} & \begin{minipage}[b]{\linewidth}\raggedright
Ratio to local FS (p50)
\end{minipage} \\
\midrule\noalign{}
\endhead
\bottomrule\noalign{}
\endlastfoot
Local FileSystem journal & \textbf{347} & 581 & 1 919 & 1.0× \\
Co-located SQL Server & 1 126 & 1 891 & 2 580 & \textbf{3.24×} \\
Co-located MySQL & 982 & 1 498 & 2 021 & \textbf{2.83×} \\
\end{longtable}
}

Durability is parameterised identically across the three backends: the
FileSystem backend invokes a per-entry flush-to-disk; SQL Server runs
with default per-commit log flush; MySQL runs with
\texttt{innodb\_flush\_log\_at\_trx\_commit=1} and
\texttt{sync\_binlog=1}. The measurement therefore compares the same
physical durability guarantee across all three --- one entry, one
durable commit.

The journal-append regime favours the substrate by ≈ 3× over both RDBMS
engines at the median, and the same ordering holds at the long tail (p95
and p99). The advantage is \emph{structural}: both RDBMS engines land in
the same ≈ 3× band, so the gap is a property of the substrate's append
being one \emph{fsync} per record versus the RDBMS being one
\emph{fsync} per redo-log entry plus the protocol overhead of a
transaction.

The corollary to E1--E4 follows. Every replay-based operation costs ≈ 3×
less per entry against the local-journal substrate than against a
co-located RDBMS in the durable-commit regime. Deployment at N entries
(E1, §5.1), replication at N entries (E2, §5.2), backup ingestion at N
entries (E3, §5.3), and catch-up after a partition at N entries (E4,
§5.4) inherit the same factor.

\paragraph{Caveats}\label{caveats-1}

Two caveats are worth naming, in the spirit of Capa 2 honesty.

First, the comparison is \emph{conservative against the substrate}. The
RDBMS engines in L6 run on ephemeral container storage --- MySQL on a
tmpfs mount, SQL Edge on a Docker named volume --- whereas the local
journal goes through the host's NVMe controller. On a production RDBMS
deployment with real disk, the RDBMS side would slow further; the 3×
ratio reported here is therefore a lower bound, not an upper one. Adding
network latency for a remote RDBMS --- a typical production topology ---
adds a flat additive term per commit, pushing the ratio further still.

Second, the comparison can be narrowed by relaxing durability on the
RDBMS side. Group commit, batched inserts, and asynchronous-durability
modes --- SQL Server's \texttt{DELAYED\_DURABILITY} and MySQL's
\texttt{innodb\_flush\_log\_at\_trx\_commit=2} --- all close part of the
gap by trading durability for throughput. These are honest engineering
choices; L6 measures the durable-commit baseline because the substrate's
local-journal regime keeps full durability and the apples-to-apples
comparison demands the same of the RDBMS side. A reader who deploys an
RDBMS-backed substrate with relaxed durability will see a smaller gap;
the structural claim --- that replay cost is bounded by append latency
--- does not change, only the constant.

\paragraph{To the next sub-section}\label{to-the-next-sub-section}

§5.5 has shown that the cost basis of the four equivalences is one
quantity --- the substrate's per-entry append latency --- and that
against a co-located RDBMS in the durable-commit regime, the
local-journal substrate runs at ≈ 3× lower latency per entry. The result
is structural rather than engine-specific: two RDBMS engines, two
implementations, the same ratio band. §5.6 closes §5 by naming two
further consequences of the substrate --- observability and debug ---
that fall out of the same property without requiring separate
apparatuses.

\subsubsection{5.6 Forensic operations as substrate
consequences}\label{forensic-operations-as-substrate-consequences}

Two further operations follow from the substrate without requiring
separate apparatuses. The substrate is a total, ordered record of every
named operation performed by the program. Conventional observability
stacks reconstruct this record indirectly by correlating logs, metrics,
and traces emitted outside the program. Under the substrate condition,
the record already exists in primary form; any observability layer
becomes a reader of the substrate rather than a reconstructor of program
history. Sagas do not appear as patterns to be inferred post-hoc from
logs; they are named at write time by the program's partition primitive
(\href{03-reactions-and-partition.md}{Paper 3}) and therefore exist in
the substrate as first-class entries.

Replay is deterministic: the same corpus, applied to the same binary,
produces the same in-memory state at the same cursor position. A runtime
issue traceable to an entry is therefore reproducible by replaying the
substrate to that entry's position; a conditional breakpoint over the
entry's identifier is trivial because the identifier is a position on
the cursor, and the cursor advances one entry at a time. Debugging
reduces from forensic analysis over logs to advancing the cursor to a
chosen position in the substrate. Both observations exist on the same
axis as §5.1--§5.5: they require no additional mechanism beyond the
substrate itself. §6 turns to the cost in operational complexity that
the canonical disciplines pay to achieve, by separate apparatuses, what
the substrate already provides.

\begin{center}\rule{0.5\linewidth}{0.5pt}\end{center}

\subsection{6. Why existing approaches duplicate work the substrate
already
does}\label{why-existing-approaches-duplicate-work-the-substrate-already-does}

The four disciplines surveyed in §2 are not unrecognised by their
respective traditions. Each has developed, over decades, a refined
apparatus to address the operational concern it names --- blue-green and
rolling-update for deployment, change-data-capture and log shipping for
replication, snapshot agents and point-in-time recovery for backup,
message-queue offsets and convergence protocols for offline operation.
The maturity of these apparatuses is not in question. What this section
examines is whether any of them dissolves the operational concern it
addresses, or whether each reconstructs, from outside the program, work
that the substrate condition of §4.1 already does.

This section is not a critique of existing practices but a structural
comparison between what those practices reconstruct and what the
substrate condition provides by construction. The four sub-sections
below examine one canonical apparatus per discipline. The pattern of
each examination is the same: name what the apparatus does, identify the
structural property it relies on, and observe that the substrate
condition provides that property by construction. The point is not that
the apparatuses are wrong but that they are responses to the absence of
the substrate condition. Once the condition holds, the work each
apparatus performs is already done by the program's own act of writing
the journal.

\subsubsection{6.1 Blue-green vs
substrate}\label{blue-green-vs-substrate}

Blue-green deployment {[}Fowler 2010{]} provisions two parallel
production environments and switches traffic between them at release
time. The apparatus has three components: a duplicate environment, a
router that decides which environment is authoritative at any moment,
and a procedure for bringing the standby environment to the state from
which it can accept traffic. The first two are operational
infrastructure; the third --- bringing the standby to a serving state
--- is what the apparatus does that the substrate condition recasts.

Under the substrate condition (§4.1), the act of bringing a process to a
state in which it can serve is the act of replaying the journal up to
the present head (E1, §4.2). The new version's instantiation differs
from the old version's continued operation only in the position of its
read cursor. The duplicate environment is therefore not a parallel
system to which state must be propagated; it is a process whose cursor
on the same corpus has not yet reached the present head. The replay that
the apparatus performs implicitly --- by re-running migration scripts,
by exporting and re-importing state, by re-priming caches --- is
performed explicitly by the substrate's reading discipline.

The structural redundancy is not that blue-green is wrong; it is that
blue-green's hardest step, the act of synchronising the standby with the
live state, is replay reconstructed manually. The reconstruction is
manual because the live state, in the canonical regime, is not the
journal of the program; it is a database whose entries record what is,
not what was said. Under the substrate condition the manual
reconstruction has nothing to do, because the corpus from which the
standby reads is the same corpus the live instance writes. The router
and the duplicate environment remain; what dissolves is the apparatus
that bridges the two.

\subsubsection{6.2 Change-data-capture vs event
streaming}\label{change-data-capture-vs-event-streaming}

Change-data-capture {[}Debezium project; Maxwell project{]} reads the
binary log of a relational store and re-emits each row-level change as
an event on an external transport. The apparatus exists because the
originating system does not emit events; it emits state changes, and CDC
observes those changes and translates them back into the event form that
downstream consumers require. The translation step is the apparatus.

The translation is informationally lossy in one direction and
inferential in the other. State changes do not carry the operation that
produced them; CDC infers, from a sequence of column-level deltas, what
the application program intended. The inference is necessarily heuristic
--- two distinct programmatic operations can produce identical
column-level deltas, and the CDC consumer cannot recover the
distinction. The apparatus is therefore not only a transport; it is a
reconstruction layer that approximates events from state, with the loss
inherent in that approximation.

Under the substrate condition the reconstruction has no work to do. The
originating system emits events --- they are the entries of the journal,
written in the language of the program (Paper 1, anti-porosity) and
reduced to compact named operations (Paper 2, separability). The wire
across which replication travels carries those entries directly (§5.2);
no inference from state, no column-level delta extraction, no schema
mapping enters the act. The CDC apparatus is structurally duplicative:
the substrate emits exactly the artifact CDC reconstructs --- minus the
reconstruction loss, and minus the apparatus that produces it.

\subsubsection{6.3 Snapshot backup vs passive
consumer}\label{snapshot-backup-vs-passive-consumer}

Snapshot backup {[}WAFL; PostgreSQL PITR; Veeam{]} captures the state of
a data store at a chosen moment and writes that capture to an archival
medium. Recovery instantiates the captured state and resumes operation
from it. The apparatus has two components: the snapshot mechanism that
produces the capture, and the recovery procedure that consumes it. Both
operate on state --- the captured artifact is a frozen image of what the
program had stored, not a record of what the program had done.

The structural property the apparatus relies on is that the program's
state, if captured atomically at a moment, suffices to characterise the
program at that moment. The property holds for programs that are
reducible to their state. It does not hold for programs that are
reducible to their history: a snapshot of the state at moment \emph{t}
does not record the operations that produced it, only their cumulative
effect, and recovery from the snapshot loses the operational record.

Under the substrate condition the program is not reducible to its state;
it is constituted by the sequence of operations in its journal (§4.1). A
consumer that holds the corpus --- without instantiating the runtime
that would replay it --- holds the program in the form in which it was
written (E3, §4.2; §5.3). The choice between holding and replaying is a
property of the destination, not of the captured artifact. The snapshot
apparatus produces a state image and discards the operational record;
the substrate emits the operational record continuously and lets the
destination decide whether and when to replay.

The redundancy is structural in two directions. First, the snapshot
apparatus is unnecessary: the substrate is already the program, copied
continuously. Second, the snapshot apparatus is also insufficient: it
captures state, which is less than the program --- recovery from a
snapshot reconstructs \emph{where} the program was, not \emph{what it
did to get there}. The passive consumer of §5.3 holds the program, and
any moment in the program's history is replayable from the corpus the
consumer holds.

\subsubsection{6.4 Message queue buffer vs persistent local
journal}\label{message-queue-buffer-vs-persistent-local-journal}

The message-queue pattern places a durable buffer between a producer and
a consumer that may not be reachable in real time. RabbitMQ's
dead-letter exchanges, Kafka's consumer offsets, and the
eventually-consistent stores derived from Dynamo {[}DeCandia et
al.~2007{]} each place an apparatus at the boundary between the system
and its counterpart. The apparatus is external to the producer; the
producer writes to its store, and a separate mechanism --- the queue,
the broker, the convergence protocol --- handles the case in which the
consumer is unreachable.

The structural property the apparatus relies on is that the producer's
store and the consumer's store are different artifacts, with a transport
between them that must be made durable. Under that structural premise,
the queue is the right apparatus: it absorbs the period during which the
consumer cannot reach the producer, and it persists the messages that
would otherwise be lost.

Under the substrate condition the producer's store and the consumer's
store are not different artifacts at the level of representation. The
producer writes to its substrate; the consumer reads from a copy of the
same substrate, or from the substrate directly. The substrate is itself
the durable artifact, and offline operation reduces to the consumer's
cursor trailing the producer's write head by an arbitrary interval (E4,
§4.2; §5.4). No queue is added between the two parties because the
substrate is already the queue --- it persists every operation in the
order it was uttered, and a consumer that has been unreachable for an
arbitrary period catches up by reading the entries that accumulated.

The redundancy here is the most structural of the four. Where the queue,
the broker, and the convergence protocol exist as boundary
infrastructure, the substrate is itself the boundary --- it is what
holds the program between writer and reader, between connected and
disconnected, between now and later. The infrastructure that the
canonical literature places between two stores is, under the substrate
condition, the property of the single store both parties share.

\subsubsection{6.5 The substrate match
table}\label{the-substrate-match-table}

The four examinations of §6.1--§6.4 share a structure. Each canonical
pattern relies on a property the substrate condition supplies by
construction; each apparatus performs work that the journal, written as
the program, already does. The pattern is summarised in the table below,
in the form of Paper 4 §6.4 {[}Paper 4 §6.4{]}: a single question ---
\emph{does the apparatus address a property the substrate already
provides?} --- answered for each pattern.

{\def\LTcaptype{none} % do not increment counter
\begin{longtable}[]{@{}
  >{\raggedright\arraybackslash}p{(\linewidth - 4\tabcolsep) * \real{0.3333}}
  >{\raggedright\arraybackslash}p{(\linewidth - 4\tabcolsep) * \real{0.3333}}
  >{\raggedright\arraybackslash}p{(\linewidth - 4\tabcolsep) * \real{0.3333}}@{}}
\toprule\noalign{}
\begin{minipage}[b]{\linewidth}\raggedright
Pattern
\end{minipage} & \begin{minipage}[b]{\linewidth}\raggedright
Property the apparatus relies on
\end{minipage} & \begin{minipage}[b]{\linewidth}\raggedright
Does the substrate condition already provide it?
\end{minipage} \\
\midrule\noalign{}
\endhead
\bottomrule\noalign{}
\endlastfoot
Blue-green deployment & The standby must be brought to the same state as
the live instance before traffic switches & Yes --- replay of the same
corpus to the present head (E1, §5.1) brings the standby to that state
by construction; no external state-propagation step is required \\
Change-data-capture & The originating store emits state changes; events
must be reconstructed from them & Yes --- the substrate emits the
operations themselves (Papers 1 and 2; §5.2); no reconstruction step is
required and no inferential loss is incurred \\
Snapshot backup & The program is characterised by its state at a moment;
recovery instantiates that state & Yes --- the program is constituted by
the sequence of operations in the journal (§4.1, §5.3); the corpus held
by a passive consumer is the program, not a state capture of it \\
Message queue buffer & A durable apparatus is required between the
producer's store and the consumer's store & Yes --- the substrate is the
durable artifact both parties share (§5.4); a consumer reading later
than the producer wrote reduces to cursor advancement over the same
corpus \\
\end{longtable}
}

The four ``Yes''s share a structure. In every case, an apparatus that
the canon developed to bridge an absence of the substrate condition
becomes structurally subsumed once the condition holds. The pattern
remains usable --- none of the four is wrong as engineering --- but the
work it performs duplicates work the substrate already does. The
maturity, sophistication, and widespread adoption of the four patterns
are not evidence that the substrate condition is unnecessary. They are
evidence that the absence of the substrate condition is costly enough to
justify four separate apparatuses to compensate for it, each developed
independently by a different tradition.

The diagnosis of §2--§6 is now complete. §7 turns from the construct to
the instantiation: a system in production whose runtime already
satisfies the substrate condition, whose four operational disciplines
fall out of the substrate without dedicated apparatus, and whose
existence is the evidence that the construct names a realisable property
and not only a conceptual one.

\begin{center}\rule{0.5\linewidth}{0.5pt}\end{center}

\subsection{7. Instantiation: realization in
Puppeteer}\label{instantiation-realization-in-puppeteer}

The construct of §§3--6 makes claims about a class of systems
characterised by the substrate condition. The present section exhibits
one such system. The voice from this point forward is concentrated and
authorial: the runtime described below is \emph{Puppeteer}, the
framework the present author maintains, and the operational decisions
named are decisions taken in deployments that run today. The framework
is presented here as an existence proof for the construct rather than as
the subject of the paper; the four equivalences hold in any runtime that
meets the substrate condition, of which Puppeteer is one.

\subsubsection{7.0 Origins}\label{origins}

\href{03-reactions-and-partition.md}{Paper 3}, §6 introduces the runtime
decomposition this section relies on: \emph{Performance} is the local
hosting of a single actor over a configured storage, \emph{Theater} is
the host for one or many Performances on a node, \emph{Ensemble} is a
pool of actors with routing, and \emph{StageManager} is the peer-to-peer
replicated state machine that coordinates topology across nodes. The
full vocabulary is presented in Paper 3 §6 and is not re-introduced
here. §7 names only the runtime surfaces on which the four equivalences
of §4.2 land: the \emph{aliveGate} (E1, §7.1), the per-event push
callback wired to a webhook delivery service (E2, §7.2), and the
\emph{Materialize} construct that declares which entries materialise in
which destinations (E3, §7.3). The offline equivalence (E4) lands on the
same write path the local-journal régime of §5.4 already exhibited, and
is referenced rather than re-developed here.

\subsubsection{7.1 Theater.aliveGate: the red-black
mechanism}\label{theater.alivegate-the-red-black-mechanism}

The handover window of §5.1 is realised in the runtime by a single
boolean predicate: the \emph{aliveGate}, a manual-reset event held on
every \texttt{Performance} instance. The gate is the locus of the
deployment apparatus, and its life-cycle is the operational realisation
of the structural intervals named in §5.1.

A \texttt{Performance} is started in one of two modes. A primary starts
with \texttt{Start()} (default \texttt{asFollower:\ false}): the runtime
hydrates the actor from its substrate, raises the
\texttt{OnFirstHydration} and \texttt{OnHydrated} callbacks, and sets
the gate. A follower starts with \texttt{Start(asFollower:\ true)}: the
runtime hydrates the actor from the same substrate, but the gate remains
reset and the public \texttt{IsAlive} predicate reports \texttt{false}.
External callers --- request dispatchers, write coordinators --- block
on \texttt{WaitUntilAlive} until the gate is set.

The handover is a single sequence over the gate. The follower invokes
\texttt{LockWhileNotSyncronized()}, which pauses writes on the live
primary and returns the EntryId at which the pause occurred. The
follower then catches up over the residual tail with
\texttt{CatchUpFromJournal(targetEntryId)}, which replays pending events
under the actor's write lock until the follower's cursor reaches the
targeted EntryId. \texttt{UnlockAndRunAlive()} flips the follower's gate
and releases the primary. From the perspective of any external caller,
the system transitions from ``primary alive'' to ``follower alive'' in a
single atomic moment --- the gate flip --- preceded by replay of the
substrate and nothing else.

The structural content of the apparatus is therefore minimal. The gate
is bookkeeping over the substrate's read cursor; the runtime contains no
separate deployment subsystem, no state-extraction layer, and no
schema-mediated handoff. The orchestration tail measured in §5.1 ---
below 30 µs at every measured N --- is the cost of the gate flip and the
residual-tail replay, not of any apparatus. Appendix A lists the exact
source locations of each named surface.

\subsubsection{7.2 Event streaming: push-to-pull
inversion}\label{event-streaming-push-to-pull-inversion}

The substrate's push-to-pull inversion (§5.2) is realised in Puppeteer
by composing three primitives that already live in the runtime. None of
them was introduced for §7.2; each existed for its own structural
reason, and together they satisfy the consumer-pull requirement.

First, the journal exposes a per-record callback.
\texttt{Diary.OnRecordWritten} is invoked by the storage layer for every
entry written, with the entry's \texttt{entryId} and its raw byte
representation. The setter accepts a single subscriber; the companion
\texttt{AddRecordWrittenCallback} chains additional subscribers without
disturbing the existing one, so that multiple downstream consumers can
attach to the same primary without the primary knowing how many.

Second, the partition primitive of
\href{03-reactions-and-partition.md}{Paper 3} exposes a \emph{Cue} mode
for Reactions. A \texttt{Cue()} Reaction runs continuously: it batches
over the substrate from the current cursor to the present head and then
enters \emph{push mode}, subscribing to the same
\texttt{OnRecordWritten} callback above. The push loop applies each
entry as it arrives, matches it against the Reaction's pattern, and
dispatches the Reaction's terminator if the pattern matches.

Third, the substrate's externalisation point is the terminator of a
\texttt{Cue} Reaction. The terminator's body --- DSL the program writes
--- is the place at which the runtime delegates to whatever external
delivery service the deployment has elected. We observe that an
off-the-shelf webhook delivery service (SVIX, in the deployments
referenced here; equivalent products would suffice) satisfies the
substrate's consumer-pull requirement: the Reaction's terminator posts
the entry to the service from a per-entry callback; the service holds
the entry durably; downstream subscribers fetch from the service at
their own pace.

No code in \texttt{Puppeteer/} or \texttt{Choreography/} references
SVIX. The runtime exposes the callback and the partition primitive; the
wiring to a particular webhook service is a property of the deployment,
not of the framework. The substrate's claim is that any service
satisfying the consumer-pull contract is sufficient; the off-the-shelf
adoption is the empirical witness to that sufficiency.

\subsubsection{7.3 Passive consumer: replicate without an actor
running}\label{passive-consumer-replicate-without-an-actor-running}

The passive consumer of §5.3 is the operational realisation that, in
earlier drafts of this paper, was named as an honest gap --- a property
derivable from the substrate condition but not yet exhibited as a system
in production. The gap closed on the day this section was drafted, by
composition of patterns the runtime already had. The composition has
since been formalised under the construct \emph{Materialize}. The
remainder of this sub-section walks the composition brick by brick, in
the spirit of the andragogical commitment of
\href{02-program-value-separability.md}{Paper 2}: a reader who has
followed §§4--6 is at risk of dismissing E3 as ``stated, not
exhibited''; the present sub-section makes the assembly visible so that
the dismissal has nothing to dismiss.

The instantiation has four bricks. The reader is asked to hold them
separately and observe how they compose.

\textbf{Brick 1 --- Same binary, alternate startup mode.} A destination
is a process running the same actor binary as the primary, started in an
alternate mode. The mode admits no external requests: the
\texttt{aliveGate} of §7.1 is not opened for the destination's lifetime.
Two such modes exist. \texttt{Performance.Start(asFollower:\ true)}
starts the actor as a \emph{follower} --- a process whose \texttt{Job()}
Reactions run against a shared substrate; the follower is a worker, not
a backup, and its Reactions contribute markers back to shared auxiliary
tables. The second mode, \texttt{/materialize}, runs no Reactions at all
--- it consumes pre-computed markers from the primary via the wire
protocol described under Brick 3. The mode is selected by the host
process at startup; the runtime is the same in either case.

\textbf{Brick 2 --- Per-event marker in the DSL.} The program declares
which entries materialise in which destinations through a marker on the
Reaction's metadata plane:

\begin{verbatim}
actor.Reactions.DefineReaction("...")
    .Job().Company().ReadForward()
    .Seek("...")
        .OnMatch("...")
    .Metadata.Materialize("DC-B");
\end{verbatim}

\texttt{Metadata.Materialize("DC-B")} records, for every entry that
matches the pattern, a row in the \texttt{EventMaterialization} table
naming the destination. The row is the program's declaration that this
entry is part of the corpus the named destination is to receive. Reading
the program reveals which entries are intended for which destinations,
in the same DSL the program is written in --- the construct is, in the
sense of \href{01-anti-porosity.md}{Paper 1}, an instance of the
homoiconic recording the substrate condition requires.

\textbf{Brick 3 --- The wire protocol between primary and destination.}
The destination, when it wakes up, asks the primary for the entries it
has not yet seen. The exchange is realised by four wire verbs on the
primary's \texttt{actor.Materialization} sub-namespace.
\texttt{ReadRecordsAfter(destination,\ fromEntryId)} returns the raw
substrate records the destination is missing --- Capa 1, the program as
the primary wrote it. \texttt{ConfirmUntil(destination,\ entryId)}
records the destination's acknowledgment on the primary side as a
Max-monotonic watermark --- the primary now knows the destination has
the corpus up to that EntryId. \texttt{ReadReactions(destination)} and
\texttt{ReadElidedRange(destination,\ from,\ to)} return Capa 2 derived
state --- the Reaction registry's atomic snapshot and the elision
markers the primary computed in the same range --- so that a destination
electing to replay can reach the same in-memory state the primary
reached without re-running the Reactions itself.

The destination side is the fluent client \texttt{MaterializeMirror}.
The contract is two-layered:

\begin{verbatim}
mirror.Sync();                       // Capa 1 — orchestrates (a) + (b).
mirror.AsProgramMirror().Sync();     // Capa 2 — orchestrates (a) + (c) + (d) + (b).
\end{verbatim}

\texttt{Sync()} fetches records and confirms the watermark;
\texttt{AsProgramMirror().Sync()} additionally fetches the Reactions
snapshot and elision markers, so that the destination's program state is
reconstructible without instantiating the Reaction matcher. The
destination decides what to do with the result --- a passive backup
persists the records and discards Capa 2; a replicated read-replica
persists both and applies the elision; a snapshot-time mirror runs Capa
1 only.

\textbf{Brick 4 --- Recovery primitive.} When a destination has been
offline for an arbitrary period --- or when an intermediate delivery
service has exhausted its retention --- the destination resumes by
issuing \texttt{Sync()} against its current watermark. The primary
serves the missing entries from its substrate; the destination's
watermark advances; the round-trip is the same as steady-state. There is
no separate recovery mode in the runtime, because there is no separate
steady-state mode either: every fetch is \texttt{Sync()} against the
watermark, and every catch-up is \texttt{Sync()} over a wider range. The
recovery primitive is the steady-state primitive.

The four bricks compose into a single operation. The program declares
which entries materialise in which destinations (Brick 2). A destination
process --- running the same binary in \texttt{/materialize} mode (Brick
1) --- fetches records from the primary's Materialization surface (Brick
3) at its own pace, confirms what it has received, and catches up after
any interruption by the same call (Brick 4). The runtime contains no
replicator, no separate backup subsystem, and no synchronisation layer
between primary and destination beyond the four wire verbs and the
per-Reaction marker. The operation falls out of the substrate; the
construct names the operation that the program already declares.

Three dimensions of the construct deserve naming, because each is
structurally present without requiring a separate apparatus.

\textbf{Heterogeneous backends.} \texttt{Diary} selects its storage
implementation polymorphically from
\texttt{(DatabaseType,\ connectionString)} at construction time. A
primary running on the local FileSystem can materialise to a destination
running on MySQL or SQL Server, and the wire protocol carries the same
raw records across the heterogeneity --- \texttt{ReadRecordsAfter}
returns Capa 1 records regardless of which backend the primary stores
them in. The four wire verbs are implemented on each of the four
backends with the same contract; L4 (§5.3) exhibits this independence
with \textbf{18 of 18 cells reaching equal in-memory state across three
backends and three journal sizes}.

\textbf{Multi-destination fanout.} \texttt{AddRecordWrittenCallback}
chains additional subscribers without displacing the existing one.
Multiple destinations may register against the same primary; the
primary's write path is unaware of how many. Each destination holds its
own watermark on the primary side; each progresses at its own rate; the
primary serves them independently. Multi-backup of one actor --- and the
per-destination retention policies that go with it --- falls out of the
chain primitive.

\textbf{Per-datacenter topology.} A datacenter is an operational
placement of \texttt{Performance} instances; it is not a construct the
actor knows about. A primary at datacenter A and a destination at
datacenter B is a topology decision realised by the operator: a local
SVIX at each site for the Cue feed of §7.2, and \texttt{/materialize}
processes at each site whose mirror clients fetch from the local SVIX or
directly from the primary. The substrate's locality property (the
journal is the actor's, not the system's; cross-ref §6.2 of
\href{03-reactions-and-partition.md}{Paper 3}) extends laterally: each
datacenter's destinations are wired to the local instance of the
delivery service, and cross-site bandwidth is consumed only by the
SVIX-to-SVIX replication the off-the-shelf service performs.

A gap is named honestly in the spirit of Capa 2. The
\texttt{Diary.OnRecordWritten} setter is wired today for primary
FileSystem and for the buffered régime; for primary MySQL or SQL Server
the setter is a silent no-op. The canonical case of the substrate (a
FileSystem-primary deployment) is unaffected, and the §7.2 push-to-pull
path runs against the FileSystem callback as designed. A future
deployment with primary MySQL would require the callback to be wired in
the MySQL storage path --- a localised patch, not a structural change.

\subsubsection{7.4 Numbers}\label{numbers}

The operational régime of the four bricks above runs under the latency
budget §5.5 already developed. The substrate cost of materialisation is
the cost of appending the program at the primary (which §5.5 measures)
plus the cost of appending at the destination (which §5.3's L4
measures). No new numbers are introduced here; the relevant ones live in
§5.1 (deployment window), §5.2 (wire density), §5.3 (destination
convergence), §5.4 (offline catch-up), and §5.5 (per-entry append
latency). Together they characterise the régime in which the four bricks
above operate; together they materialise the construct in
production-grade numbers.

\begin{center}\rule{0.5\linewidth}{0.5pt}\end{center}

\subsection{8. Relation to previous work in this paper
series}\label{relation-to-previous-work-in-this-paper-series}

The substrate condition stated in §4.1 --- that the journal is
homoiconic, dense, and made of compact named operations --- is not
asserted ex nihilo by the present paper. Each of its components is the
conclusion of prior structural analysis in this series, and each entered
§4.1 as a precondition rather than as a claim of this paper. Without
those preconditions the four equivalences of §4.2 would not follow; with
them, the four equivalences are corollaries of the substrate's
properties under variations of when and where the read happens. §8 makes
the chain of dependence explicit.

\href{01-anti-porosity.md}{Paper 1} introduces \emph{anti-porosity} as a
design principle for the boundary between domain code and persisted
form. The journal records what the program said, not the data structures
the program manipulated; entries are sentences in the same language the
program is written in, and every effect that crosses the boundary is
named in that language rather than summarised as a state delta.
Anti-porosity is the precondition under which the journal can be the
program at all: a porous journal records states whose operational origin
is lost, and the substrate condition reduces to recording a derivative
of the program rather than the program itself. §4.1 takes anti-porosity
as given; without it, E2 (replication as sharing history) and E3 (backup
as copying the program) collapse into the canonical regimes of §2.

\href{02-program-value-separability.md}{Paper 2} introduces
\emph{separability} as the structural property that makes the journal
entry of a parametric verb a compact reference plus its arguments. The
verb's body is not present in the entry; it lives once, named, in the
corpus's vocabulary. Separability is the precondition under which the
wire of §5.2 carries operations at a density that does not scale with
the operation's internal complexity, and it is the precondition under
which the latency budget of §5.5 holds --- replay is linear in
compact-entry count, not in expanded state, only because separability
has reduced the entry to its reference form. §5.5's three-times
advantage against co-located RDBMS engines is, in part, the operational
expression of the structural property Paper 2 names.

\href{03-reactions-and-partition.md}{Paper 3} introduces \emph{the
partition} between immediate work, done at the actor's writer-lock
cursor, and deferred work, done by Reactions that read the journal
forward and produce their own derived entries. Paper 3 §6 establishes
the runtime decomposition --- \emph{Performance}, \emph{Theater},
\emph{Ensemble} --- that frames the operational realisation of §7. The
mechanism that admits a new instance as authoritative (§5.1's gate flip)
is named in Paper 3 claim 8 as a property derivable from the partition;
the same paper names the cross-actor extension as a forward reference
closed by Paper 4. §5.2's symmetric-consumer property --- that two
instances of the same binary, reading the same corpus, produce identical
Reaction output (L3) --- depends on the \texttt{expose} discipline Paper
3 §5 establishes: a Reaction's guard reads only data the program has
explicitly exposed, so whatever the remote consumer needs to evaluate
the guard travels in the corpus the program writes.

\href{04-cross-actor-continuity.md}{Paper 4} introduces \emph{tell} as
the primitive under which cross-actor causation is recorded as program
in each participant's local journal. Paper 4 claim 10 names
cross-datacenter replication as the canonical case in which tell's
program-level recording survives where the conventional apparatus (CDC,
log shipping) reconstructs state and loses the causal chain. §5.2 of the
present paper completes the picture from the substrate side: the journal
record is what the consumer reads, and the consumer reconstructs not
only the local actor's state but, when tell entries are present, the
cross-actor causal chain those entries record. The substrate condition
makes Paper 4's cross-DC property hold without an external mechanism
that observes the producer's runtime.

The dependence is summarised in the table below.

{\def\LTcaptype{none} % do not increment counter
\begin{longtable}[]{@{}
  >{\raggedright\arraybackslash}p{(\linewidth - 6\tabcolsep) * \real{0.2500}}
  >{\raggedright\arraybackslash}p{(\linewidth - 6\tabcolsep) * \real{0.2500}}
  >{\raggedright\arraybackslash}p{(\linewidth - 6\tabcolsep) * \real{0.2500}}
  >{\raggedright\arraybackslash}p{(\linewidth - 6\tabcolsep) * \real{0.2500}}@{}}
\toprule\noalign{}
\begin{minipage}[b]{\linewidth}\raggedright
Paper
\end{minipage} & \begin{minipage}[b]{\linewidth}\raggedright
What it establishes
\end{minipage} & \begin{minipage}[b]{\linewidth}\raggedright
What §4.1 takes as given
\end{minipage} & \begin{minipage}[b]{\linewidth}\raggedright
Where it is load-bearing in this paper
\end{minipage} \\
\midrule\noalign{}
\endhead
\bottomrule\noalign{}
\endlastfoot
\href{01-anti-porosity.md}{Paper 1} & Anti-porosity: the journal records
operations, not state & The journal is homoiconic and dense & E2 wire
density (§5.2); E3 the corpus is the program (§5.3); §6.3 snapshot is
informationally less \\
\href{02-program-value-separability.md}{Paper 2} & Separability: the
entry of a parametric verb is \texttt{ActionId} + arguments & The
journal entry is compact, named, parameter-referenced & Wire density at
production scale (§5.2; L2); replay rate linear in compact count, not
state (§5.1; L1); latency budget (§5.5; L6) \\
\href{03-reactions-and-partition.md}{Paper 3} & Partition: immediate vs
deferred; Reactions; runtime decomposition (\emph{Performance},
\emph{Theater}, \emph{Ensemble}) & Reactions and the partition
primitive; the runtime apparatus referenced in §7 & Gate flip apparatus
(§5.1; §7.1); symmetric consumer via \emph{expose} (§5.2; L3); apparatus
vocabulary (§7.0) \\
\href{04-cross-actor-continuity.md}{Paper 4} & \emph{tell} primitive:
cross-actor causation as program in sender's journal & The cross-actor
recording the substrate carries across sites & Cross-DC causal chain
travels with replication (§5.2; cross-ref claim 10) \\
\end{longtable}
}

The four prior papers can be read as the structural derivation that the
present paper takes as substrate. Each named one property that the
journal had to possess for the program to live in it; the present paper
takes the four properties as given and reads off the four operational
equivalences that follow. Where each prior paper went from one
structural commitment to its consequences, the present paper goes from
four structural commitments to a single primitive --- replay --- under
which the four canonical operational disciplines collapse. The relation
is therefore not parallel: §4.1's substrate is what the prior four
established, and §4.2's theorem is what follows once it is in hand.

\begin{center}\rule{0.5\linewidth}{0.5pt}\end{center}

\subsection{9. Counter-arguments}\label{counter-arguments}

The substrate condition and its four equivalences invert vocabulary the
canonical literature has used productively for decades. A reader
familiar with that literature is likely to raise objections that are not
addressed by §§2--8. Four are taken up here. The treatment of each
follows the same shape: name the valid aspect of the objection, then
identify the structural premise on which the invalid aspect depends and
observe that the substrate condition removes that premise. The
refutations are made from the construct, not from any framework that
instantiates it.

\subsubsection{9.1 ``Loading a large journal takes
time''}\label{loading-a-large-journal-takes-time}

\textbf{Canonical assumption.} Bringing a system to a serving state
requires reconstructing its persisted history end-to-end, and the cost
of that reconstruction grows without bound with the system's operational
age.

The objection. A journal that accumulates entries over the operational
life of a program grows without bound. Bringing a new instance to a
serving state by replaying the journal therefore takes longer the older
the program becomes. A canonical anecdote names a relational system
whose initial load --- from a binary log accumulated over a year ---
required roughly an hour before the system could accept traffic. If
deployment is replay (§5.1), then deployment is bounded below by that
load time, and the substrate's account of deployment makes the deploy
window worse, not better.

What is valid. Replay is linear in the entries it reads. A journal twice
as long takes twice as long to replay. The objection correctly observes
that the absolute deploy time grows with the program's history under any
substrate-based account of deployment.

What does not follow. The hour-long load named in the anecdote is not
replay over the substrate condition; it is reconstruction of state from
a log whose entries are state-change deltas. The two régimes have
different cost basis. Under the substrate condition, the entry is a
compact named operation --- \texttt{ActionId} plus arguments --- and
replay is the application of those operations in order (Papers 1 and 2).
The per-entry cost is the in-memory apply of one operation, not the
disk-bound write or schema-mediated commit that state-change replay
incurs.

§5.1's measurement materialises the structural distinction. L1 reports a
sustained replay rate of \textbf{451 thousand entries per second} at the
median (N = 100 000) and \textbf{416 thousand entries per second} at N =
1 000 000 --- the rate is bounded by neither N nor the per-entry apply
cost, and the orchestration tail around bulk replay is below 30 µs at
every measured N. A journal of ten million compact entries replays in
the order of twenty seconds at this rate, not an hour. The objection
generalises a measurement made in one régime --- state replay over a
binary log --- to a régime in which the entries are operations, and the
generalisation does not hold. What grows linearly is the number of
operations the program has uttered; what does not grow is the cost of
uttering each one again in memory.

One canonical case lies outside this argument and is preserved as an
honest limit. Where the cost-dominant operation at a single entry is
itself expensive --- a network round-trip to an external service, a
complex domain calculation --- replay carries that cost per entry just
as live operation did. The substrate does not make any one operation
cheaper; it makes the bookkeeping around the operation negligible.

\subsubsection{9.2 ``Change-data-capture already solves
replication''}\label{change-data-capture-already-solves-replication}

\textbf{Canonical assumption.} Cross-datacenter replication is properly
addressed by reconstructing events from the originating store's
state-change deltas, mediated by a CDC pipeline between primary and
replica.

The objection. Cross-datacenter replication is a mature engineering
concern. Change-data-capture pipelines, log shipping between database
engines, and managed replication services route writes from primary to
replica with bounded lag and well-understood failure modes. The
substrate's account of replication (§5.2) appears to re-derive a problem
that the canon has solved repeatedly. What additional structural claim
does the substrate condition support that the existing apparatus does
not?

What is valid. CDC pipelines work. Production systems replicate across
datacenters with CDC every day; the apparatus is reliable,
vendor-supported, and well-tooled. As an engineering solution, it solves
the operational concern of moving state between sites.

What does not follow. The structural premise of CDC is that the
originating store records state changes, not operations, and that events
must be reconstructed from those state changes by inferential reading of
the binary log (§2.2, §6.2). The reconstruction is lossy: two distinct
programmatic operations that produce identical state changes are
indistinguishable at the CDC layer. The wire across which CDC traffic
travels carries a derived artifact --- column-level deltas --- at a
width that scales with the size of the affected rows rather than with
the operation that caused them.

The substrate condition replaces the reconstruction with direct
emission. The journal entry is the program's operation, in the language
of the program (Paper 1), reduced to a reference to a parametric verb
plus its arguments (Paper 2). The wire carries the operation as uttered,
not the state changes it produced. §5.2's measurement (L2) reports a
compact wire of \textbf{36 to 99.7 bytes per event} depending on verb
tier, against literal-form payloads of \textbf{121 to 913.7 bytes}; the
literal-to-compact ratio reaches \textbf{9.2× at the production-shaped
tier and 20× at full production scale} (Paper 2 Lab 4). Applying
\texttt{gzip} to the literal form closes the gap only to \textbf{3.9× at
tier 3} --- the structural saving is not algorithmic but
representational, and a compressor operating after the fact cannot
recover what naming achieved by construction.

The objection is therefore precise in scope. CDC solves replication as
state propagation. The substrate condition does not deny that solution;
it shows that under the construct's conditions the problem CDC solves
does not arise, and the wire carries a different artifact. The two are
not in competition for the same engineering decision; they are at
different layers of the design.

\subsubsection{9.3 ``Snapshot backup is
simpler''}\label{snapshot-backup-is-simpler}

\textbf{Canonical assumption.} Backup is a captured image of state at a
chosen moment, archived by an external apparatus and re-instantiated at
recovery time.

The objection. Backup as a periodic state snapshot is a forty-year-old
discipline with well-understood operational properties: recovery time
objective (RTO), recovery point objective (RPO), storage tiering,
lifecycle policies. The substrate's account of backup (§5.3) requires a
continuously consuming destination --- a process that holds the corpus
as it grows --- which appears operationally more complex than periodic
state capture. Why prefer continuous to periodic?

What is valid. Periodic snapshots have an operational vocabulary the
field knows. The discipline is mature, vendor-supported, and integrates
with object-store lifecycle infrastructure. As an engineering choice for
systems whose program is reducible to state, the snapshot apparatus is a
sound default.

What does not follow. The structural premise of snapshot backup is that
the program is reducible to its state --- that a captured image at
moment \emph{t} suffices to characterise the program at that moment
(§6.3). The premise holds for programs whose journal is, in effect, a
derivative of state changes. It does not hold for programs whose journal
is the program itself.

Under the substrate condition, the artifact that characterises the
program is the journal, not any state derived from it. A snapshot at
moment \emph{t} captures \emph{where} the program had arrived; the
corpus to \emph{t} captures \emph{how it got there}. The two are not
equivalent informational artifacts: from the snapshot, the operational
history is lost; from the corpus, the snapshot at any past moment is
recoverable by replay. The substrate-based account of backup is
therefore not simply a different operational pattern with the same
content; it is the strictly more informative choice when the construct's
conditions hold.

The cost is honest. A continuously consuming destination, even one whose
role is to hold the corpus and never replay it, must be reachable, must
persist what arrives, and must catch up when it falls behind. §5.3's
measurement (L4) reports steady-state convergence of \textbf{≈ 2 419
events/sec on FileSystem, ≈ 1 004 events/sec on MySQL, and ≈ 764
events/sec on SQL Server}, with per-instance verification reaching equal
in-memory state in \textbf{18 of 18 cells across three backends and
three journal sizes}. The numbers establish that the operational
obligation is bounded --- a destination that ingests at the rates above
is not a degenerate case --- and that the corpus arrives identically
across heterogeneous backends. The objection's premise that simplicity
favours snapshot is honest engineering; the construct's reply is that
the simpler artifact carries less information, and the choice between
holding state and holding program is the choice the substrate condition
makes available.

\subsubsection{9.4 ``Append does not scale beyond N
writes/sec''}\label{append-does-not-scale-beyond-n-writessec}

\textbf{Canonical assumption.} A journaled append is bounded by the
per-store write throughput of the underlying medium, and that ceiling
constrains any substrate-based account to per-actor write rates the
medium can absorb.

The objection. Every persistent store has a write-throughput ceiling. A
substrate that records every operation as a journal append inherits
whatever ceiling the underlying store imposes; production systems that
exceed that ceiling are constrained either to batched-write modes or to
relaxed-durability commits. The substrate condition therefore does not
generalise to high-throughput workloads, and the four equivalences hold
only for systems that fit within the per-actor write rate the journal
can absorb.

What is valid. A local-journal substrate, like any persistent store, has
a per-entry append latency that bounds its write throughput. §5.5's
measurement (L6) reports \textbf{347 µs p50 on the local FileSystem
backend} under one-entry one-durable-commit semantics, corresponding to
roughly three thousand entries per second of single-writer append rate.
A workload that exceeds that rate on a single actor cannot be absorbed
by single-writer append at the same durability guarantee.

What does not follow. The single-actor ceiling is not a ceiling on the
construct; it is a property of the configuration in which one actor
handles one event stream. The substrate condition is preserved under
sharding (one actor per shard, each with its own journal --- the
per-actor invariant of \href{03-reactions-and-partition.md}{Paper 3},
§6.2 and §6.8) and under local buffering with asynchronous replication
to a canonical backend (§5.4's offline equivalence applies to the
steady-state régime, not only to outage absorption). Neither sharding
nor buffering is foreign to the construct; both fall out of the
substrate's locality property --- the journal is the actor's, not the
system's.

The honest scope is therefore narrower than the objection suggests. The
construct does not claim that a single actor with one journal absorbs
unbounded write rates; it claims that within the rate the journal can
absorb at a given durability guarantee, the four equivalences hold.
§5.5's table records that the local-journal régime favours the substrate
by \textbf{≈ 3× over both co-located RDBMS engines at p50 and at the
long tail}, and that the advantage is structural (both RDBMS engines
land in the same band). The trade-off is not journal-vs-no-journal; it
is local-append-vs-network-commit, and the substrate's ceiling is the
higher of the two before any sharding is considered.

Relaxing durability changes the constant. Group commit, batched inserts,
and asynchronous-durability modes on the RDBMS side close part of the
gap by trading durability for throughput; the substrate's local-append
rate can be raised similarly by group-flush semantics. §5.5's caveats
name both choices as honest engineering. What the comparison preserves
under all durability régimes is the structural claim: replay-based
operations inherit append latency as their dominant cost term, and a
local append against a substrate is, in every measured régime, at least
as fast as a comparable commit against a co-located transactional store.

\begin{center}\rule{0.5\linewidth}{0.5pt}\end{center}

\subsection{10. Conclusion}\label{conclusion}

The canonical literature on distributed systems has long treated
deployment, cross-datacenter replication, backup, and offline operation
as four operational disciplines, each supported by its own apparatus.
This treatment proved productive: blue-green environments,
change-data-capture pipelines, snapshot agents, and message-queue
buffers matured into reliable engineering practice. From that
productivity grew the fragmentation analysed in §2 and §6 --- four
parallel literatures, four parallel apparatuses, and four parallel cost
models, borne independently by systems that must be zero-downtime,
cross-datacenter, recoverable, and tolerant of disconnection.

This paper observes that the fragmentation is not entailed by the
operations themselves. It follows from the assumption named in §3: that
operational concerns are addressed by apparatus around the program
rather than by the program itself. Under the substrate condition
established by prior papers in this series, that assumption no longer
holds. The journal is the program written out over time, and four
equivalences follow from reading it under variations of when and where
(§4.2): deployment is replay at the present head; replication is the
same replay at another site; backup is the corpus held without replay;
offline operation is replay after a delay.

Viewed from this condition, the apparatuses surveyed in §6 appear as
mechanisms that reconstruct, from outside the program, work that the
substrate already performs by construction. Blue-green deployment
reconstructs replay manually; change-data-capture reconstructs
operations from state; snapshot backup captures a derivative of the
program's history; message-queue buffers introduce a durable artefact to
absorb disconnection that the substrate's own persistence already
absorbs.

The contribution of this paper is conceptual. The instantiation in
production serves as an existence proof; the measurements in §5 exhibit
the régime in which the four equivalences are observable. One artefact
--- the substrate --- supports four operational readings without
requiring separate mechanisms.

For a journaled program, deployment is replay, replication is sharing
history, backup is copying the program, and offline operation is delayed
replay. The four operational disciplines named in §2 remain meaningful;
what changes is the structural locus of operational concerns: from
externally-added apparatus to internal readings of the program's
substrate.

\begin{center}\rule{0.5\linewidth}{0.5pt}\end{center}

\subsection{Appendix A. Source-code
verification}\label{appendix-a.-source-code-verification}

The constructs named in §7 are publicly available in the Puppeteer
codebase. The line references below are pinned against the master branch
as of the drafting of this section. Three tables organise the surface by
sub-section: the \emph{aliveGate} mechanism (§7.1), the push-to-pull
primitives (§7.2), and the \emph{Materialize} construct (§7.3).

\subsubsection{§7.1 --- Theater.aliveGate}\label{theater.alivegate}

{\def\LTcaptype{none} % do not increment counter
\begin{longtable}[]{@{}
  >{\raggedright\arraybackslash}p{(\linewidth - 6\tabcolsep) * \real{0.2500}}
  >{\raggedright\arraybackslash}p{(\linewidth - 6\tabcolsep) * \real{0.2500}}
  >{\raggedright\arraybackslash}p{(\linewidth - 6\tabcolsep) * \real{0.2500}}
  >{\raggedright\arraybackslash}p{(\linewidth - 6\tabcolsep) * \real{0.2500}}@{}}
\toprule\noalign{}
\begin{minipage}[b]{\linewidth}\raggedright
Surface
\end{minipage} & \begin{minipage}[b]{\linewidth}\raggedright
File
\end{minipage} & \begin{minipage}[b]{\linewidth}\raggedright
Line(s)
\end{minipage} & \begin{minipage}[b]{\linewidth}\raggedright
What it shows
\end{minipage} \\
\midrule\noalign{}
\endhead
\bottomrule\noalign{}
\endlastfoot
\texttt{aliveGate} manual-reset event &
\texttt{Choreography/Theater/Performance.cs} & 34 & Reset while follower
or during handover; Set when alive \\
\texttt{Start(bool\ asFollower\ =\ false)} &
\texttt{Choreography/Theater/Performance.cs} & 100 & Two-mode start;
sets gate iff primary \\
\texttt{WaitUntilAlive(CancellationToken)} &
\texttt{Choreography/Theater/Performance.cs} & 136 & External callers
block here \\
\texttt{IsAlive} predicate &
\texttt{Choreography/Theater/Performance.cs} & 167 & False while
follower; true after handover \\
\texttt{LockWhileNotSyncronized()} &
\texttt{Choreography/Theater/Performance.cs} & 177 & Pause writes on
primary; return EntryId \\
\texttt{UnlockAndRunAlive()} &
\texttt{Choreography/Theater/Performance.cs} & 183 & Flip gate; release
primary \\
\texttt{CatchUpFromJournal(long\ targetEntryId)} &
\texttt{Puppeteer/EventSourcing/ActorHandler.cs} & 2159 & Replay
residual tail under write lock \\
\texttt{RehydrateFromEvent} &
\texttt{Puppeteer/EventSourcing/DB/Diary.cs} & 228 & Diary-side bulk
replay primitive \\
\end{longtable}
}

\subsubsection{§7.2 --- Push-to-pull
primitives}\label{push-to-pull-primitives}

{\def\LTcaptype{none} % do not increment counter
\begin{longtable}[]{@{}
  >{\raggedright\arraybackslash}p{(\linewidth - 6\tabcolsep) * \real{0.2500}}
  >{\raggedright\arraybackslash}p{(\linewidth - 6\tabcolsep) * \real{0.2500}}
  >{\raggedright\arraybackslash}p{(\linewidth - 6\tabcolsep) * \real{0.2500}}
  >{\raggedright\arraybackslash}p{(\linewidth - 6\tabcolsep) * \real{0.2500}}@{}}
\toprule\noalign{}
\begin{minipage}[b]{\linewidth}\raggedright
Surface
\end{minipage} & \begin{minipage}[b]{\linewidth}\raggedright
File
\end{minipage} & \begin{minipage}[b]{\linewidth}\raggedright
Line(s)
\end{minipage} & \begin{minipage}[b]{\linewidth}\raggedright
What it shows
\end{minipage} \\
\midrule\noalign{}
\endhead
\bottomrule\noalign{}
\endlastfoot
\texttt{Diary.OnRecordWritten} setter &
\texttt{Puppeteer/EventSourcing/DB/Diary.cs} & 145 & Per-record callback
(single subscriber) \\
\texttt{AddRecordWrittenCallback} &
\texttt{Puppeteer/EventSourcing/DB/Diary.cs} & 160 & Chain N subscribers
without displacing existing \\
\texttt{ReactionMode.Cue} &
\texttt{Puppeteer/EventSourcing/Follower/Reaction.cs} & 19 & Continuous
push-mode Reaction \\
Cue batch → push loop wiring &
\texttt{Puppeteer/EventSourcing/Follower/Reaction.cs} & 563-579 &
Catch-up batch then subscribe to callback \\
SVIX integration & --- & --- & Not present in \texttt{Puppeteer/} or
\texttt{Choreography/}; wired at the per-API host process \\
\end{longtable}
}

\subsubsection{§7.3 --- Materialize
construct}\label{materialize-construct}

{\def\LTcaptype{none} % do not increment counter
\begin{longtable}[]{@{}
  >{\raggedright\arraybackslash}p{(\linewidth - 6\tabcolsep) * \real{0.2500}}
  >{\raggedright\arraybackslash}p{(\linewidth - 6\tabcolsep) * \real{0.2500}}
  >{\raggedright\arraybackslash}p{(\linewidth - 6\tabcolsep) * \real{0.2500}}
  >{\raggedright\arraybackslash}p{(\linewidth - 6\tabcolsep) * \real{0.2500}}@{}}
\toprule\noalign{}
\begin{minipage}[b]{\linewidth}\raggedright
Surface
\end{minipage} & \begin{minipage}[b]{\linewidth}\raggedright
File
\end{minipage} & \begin{minipage}[b]{\linewidth}\raggedright
Line(s)
\end{minipage} & \begin{minipage}[b]{\linewidth}\raggedright
What it shows
\end{minipage} \\
\midrule\noalign{}
\endhead
\bottomrule\noalign{}
\endlastfoot
\texttt{Performance.Start(asFollower:\ true)} &
\texttt{Choreography/Theater/Performance.cs} & 100, 118 & Alternate
startup mode; \texttt{SuppressReactionJournaling} flag \\
\texttt{Reaction.MaterializeFromMetadata(destination)} &
\texttt{Puppeteer/EventSourcing/Follower/Reaction.cs} & 669 & DSL marker
plumbing: records destination on action \\
\texttt{actor.Materialization} sub-namespace &
\texttt{Puppeteer/Materialization.cs} & (whole file) & Register /
Deregister / List / ReadRecordsAfter / ConfirmUntil / ReadReactions /
ReadElidedRange \\
\texttt{MaterializeMirror} fluent client &
\texttt{Puppeteer/MaterializeMirror.cs} & (whole file) & \texttt{Sync()}
(Capa 1); \texttt{AsProgramMirror().Sync()} (Capa 2) \\
\texttt{MirrorSyncResult} immutable struct &
\texttt{Puppeteer/MaterializeMirror.cs} & 170-202 & Records /
ReactionsSnapshot / ElisionMarkers + watermark advance \\
Diary backend polymorphism &
\texttt{Puppeteer/EventSourcing/DB/Diary.cs} & 53-78 &
\texttt{(DatabaseType,\ connectionString)} → InMemory / FileSystem /
MySQL / SQLServer \\
Local-buffer (offline régime) &
\texttt{Puppeteer/EventSourcing/DB/Diary.cs} & 28, 34, 80-135 &
\texttt{IsBuffered}; \texttt{ReplicationAgent} async drain to canonical
storage \\
\end{longtable}
}

\begin{center}\rule{0.5\linewidth}{0.5pt}\end{center}

\subsection{Code provenance}\label{code-provenance}

Source-code references in this paper resolve against the public
Puppeteer repository at commit
\href{https://github.com/alvaroNCubo/puppeteer/tree/2f31f9674a5de816bdf1bf9d8360ff218a02e4da}{\texttt{2f31f96}}
(2026-05-18). The snapshot is archived in Software Heritage under the
following persistent identifier:

\begin{verbatim}
swh:1:dir:10e7e6bad7eb77b6c2e406762026177f95c3ae92;
  origin=https://github.com/alvaroNCubo/puppeteer;
  anchor=swh:1:rev:2f31f9674a5de816bdf1bf9d8360ff218a02e4da
\end{verbatim}

Inline references of the form \texttt{file.cs:NN} (e.g.,
\texttt{ActorHandler.cs:38}) resolve against this snapshot. A reader can
construct a per-file SWHID by adding the qualifiers
\texttt{;path=\textless{}path\textgreater{};lines=\textless{}NN\textgreater{}}
to the directory SWHID above. Future commits to the repository may
renumber lines; the SWHID preserves the cited state independently of any
future change to the repository or its hosting.

\subsection{Acknowledgments}\label{acknowledgments}

The author used large language models (including Claude and ChatGPT) as
editorial assistants for language refinement, structural feedback, and
literature navigation. All original ideas, terminology, theoretical
constructs, and technical content presented in this work are solely the
author's.

\begin{center}\rule{0.5\linewidth}{0.5pt}\end{center}

\subsection{Appendix B. Bibliography}\label{appendix-b.-bibliography}

\emph{Online sources accessed at the moment of drafting; dates updated
at release.}

\begin{itemize}
\tightlist
\item
  Ambler, S. W., \& Sadalage, P. J. (2006). \emph{Refactoring Databases:
  Evolutionary Database Design.} Addison-Wesley. ISBN 978-0-321-29353-5.
\item
  AWS documentation. \emph{Managing your storage lifecycle.}
  https://docs.aws.amazon.com/AmazonS3/latest/userguide/object-lifecycle-mgmt.html
  (accessed 2026-05-14).
\item
  Azure documentation. \emph{Overview of operational backup for Azure
  Blobs.}
  https://learn.microsoft.com/en-us/azure/backup/blob-backup-overview
  (accessed 2026-05-14).
\item
  Bacula project. https://www.bacula.org/ (accessed 2026-05-14).
\item
  Burns, B., Grant, B., Oppenheimer, D., Brewer, E., \& Wilkes, J.
  (2016). \emph{Borg, Omega, and Kubernetes.} Communications of the ACM,
  59(5), 50--57.
\item
  Debezium project. https://debezium.io/ (accessed 2026-05-14).
\item
  DeCandia, G., Hastorun, D., Jampani, M., Kakulapati, G., Lakshman, A.,
  Pilchin, A., Sivasubramanian, S., Vosshall, P., \& Vogels, W. (2007).
  \emph{Dynamo: Amazon's Highly Available Key-value Store.} Proceedings
  of the 21st ACM SIGOPS Symposium on Operating Systems Principles (SOSP
  '07), 205--220.
\item
  Flyway project. https://flywaydb.org/ (accessed 2026-05-14).
\item
  Fowler, M. (2010). \emph{BlueGreenDeployment.} martinfowler.com/bliki.
  https://martinfowler.com/bliki/BlueGreenDeployment.html (accessed
  2026-05-14).
\item
  Hevner, A. R., March, S. T., Park, J., \& Ram, S. (2004). \emph{Design
  science in information systems research.} MIS Quarterly, 28(1),
  75--105.
\item
  Hitz, D., Lau, J., \& Malcolm, M. (1994). \emph{File System Design for
  an NFS File Server Appliance.} Proceedings of the USENIX Winter
  Technical Conference, San Francisco, January 1994.
\item
  Humble, J., \& Farley, D. (2010). \emph{Continuous Delivery: Reliable
  Software Releases through Build, Test, and Deployment Automation.}
  Addison-Wesley. ISBN 978-0-321-60191-9.
\item
  Kreps, J., Narkhede, N., \& Rao, J. (2011). \emph{Kafka: a Distributed
  Messaging System for Log Processing.} Proceedings of the NetDB
  Workshop, 6th International Workshop on Networking Meets Databases
  (co-located with SIGMOD 2011).
\item
  Kubernetes documentation. \emph{Deployments --- Rolling Update
  Deployment Strategy.}
  https://kubernetes.io/docs/concepts/workloads/controllers/deployment/\#rolling-update-deployment
  (accessed 2026-05-14).
\item
  Lakshman, A., \& Malik, P. (2010). \emph{Cassandra: A Decentralized
  Structured Storage System.} ACM SIGOPS Operating Systems Review,
  44(2), 35--40.
\item
  Liquibase project. https://www.liquibase.org/ (accessed 2026-05-14).
\item
  Maxwell project. \emph{Maxwell's Daemon.} https://maxwells-daemon.io/
  (accessed 2026-05-14).
\item
  MySQL Reference Manual. \emph{Chapter 17: Replication.}
  https://dev.mysql.com/doc/refman/8.0/en/replication.html (accessed
  2026-05-14).
\item
  PostgreSQL documentation. \emph{Continuous Archiving and Point-in-Time
  Recovery (PITR).}
  https://www.postgresql.org/docs/current/continuous-archiving.html
  (accessed 2026-05-14).
\item
  RabbitMQ documentation. \emph{Dead Letter Exchanges.}
  https://www.rabbitmq.com/dlx.html (accessed 2026-05-14).
\item
  Redis documentation. \emph{Redis replication.}
  https://redis.io/docs/management/replication/ (accessed 2026-05-14).
\item
  Sato, D. (2014). \emph{CanaryRelease.} martinfowler.com/bliki.
  https://martinfowler.com/bliki/CanaryRelease.html (accessed
  2026-05-14).
\item
  Veeam documentation. \emph{Veeam Backup \& Replication User Guide.}
  https://helpcenter.veeam.com/docs/backup/vsphere/overview.html
  (accessed 2026-05-14).
\end{itemize}

\begin{center}\rule{0.5\linewidth}{0.5pt}\end{center}

\begin{itemize}
\tightlist
\item
  Rivera, A. (2026). \emph{Anti-porous architecture: a unified design
  principle for CQRS + Actor + Event-Sourcing systems.} Paper 1 of this
  series. \url{01-anti-porosity.md}
\item
  Rivera, A. (2026). \emph{Program-value separability: the structural
  precondition for compilation, caching, and dense journaling in a DSL
  runtime.} Paper 2 of this series.
  \url{02-program-value-separability.md}
\item
  Rivera, A. (2026). \emph{Reactions and the partition: opt-in eventual
  consistency in actor-native systems.} Paper 3 of this series.
  \url{03-reactions-and-partition.md}
\item
  Rivera, A. (2026). \emph{Preserving semantic continuity across actors:
  a tell-based approach without orchestration.} Paper 4 of this series.
  \url{04-cross-actor-continuity.md}
\end{itemize}

\end{document}
